\documentclass[../main.tex]{subfiles}

\graphicspath{{\subfix{../images/}}}

\begin{document}

\clearpage
\thispagestyle{empty}
\vspace*{7cm}
\section{Measurement uncertainty estimate}
\vspace*{1.5cm}

\subsection{Introduction}

Estimating measurement uncertainty is one of the most important steps in performing a measurement, since uncertainty—which must always be reported—reflects the degree of uncertainty in our knowledge of the measurand. 

An underestimation of uncertainty leads to attributing greater significance to a measurement than it actually has, resulting in technical and/or legal problems. On the other hand, an overestimation undervalues the quality of the measurement, which was obtained through the use of expensive equipment and complex methods. 

It is important to emphasize that the uncertainty of a measurement cannot be reduced to zero because: 

\begin{enumerate}
	\item[\textbf{-}] The equipment and samples used to perform the measurement are not ideal;
	\item[\textbf{-}] The interaction between the equipment and the system under measurement alters the state of the system itself, so the resulting measurement will in any case differ from that obtained under ideal conditions;
	\item[\textbf{-}] The measurements of the quantities that define the state of the system under measurement and of the influencing quantities are characterized by non-zero uncertainty;
	\item[\textbf{-}] The definition of the quantity being measured, however complex, cannot be exhaustive.
\end{enumerate}

\subsection{Classification of measurement methods}


An initial classification of measurement methods, based on the number of readings taken to assign a value to the measured parameter, allows us to distinguish between:\\

\textbf{single-reading} measurement methods, in which the measurement is assigned following the execution of a single reading on each of the instruments involved;\\

measurement methods with \textbf{repeated readings}, which involve taking multiple readings of each quantity under nominally identical conditions, and then assigning the measurement as the result of a statistical analysis of the data set thus obtained.\\

Another way to classify measurement methods is one that takes into account the operational procedures for assigning a measurement to a parameter, which allows us to identify the following two categories:\\

\textbf{direct measurement methods} are those that allow a measurement to be assigned to a parameter based on  a reading, or a series of readings, from an instrument without needing to know other parameters of the measured system, except for the value of any samples, the influencing quantities, and the quantities explicitly referenced in the definition of the measurand; \\

\textbf{indirect measurement methods} are those in which the measurement of a parameter is assigned as the result of a calculation involving the values of other parameters measured directly.

\subsubsection{Direct measurement methods}
\label{221}

Direct measurement methods are based on comparing the quantity being measured with a quantity of the same type generated by (or stored in) a local sample.

The most common direct measurements are those performed using the a direct reading method, i.e., assigning the measurement $m$ based on the indication provided by a device to which the measurand is applied via the following relationship:


\begin{equation}
	m = f_t(l)
\end{equation}

where $f_t$ is the calibration function stored in the instrument being used.
This measurement method therefore requires that the instrument be calibrated before use, in order to establish the correspondence between the readings provided by the instrument and the measurements of the quantities applied to its input.

The main sources of uncertainty in the measurements obtained using this technique are:

\begin{enumerate}
	\item[\textbf{-}] Instrumental uncertainty, which is specified by the manufacturer and usually depends on the values of significant influencing factors and the time elapsed since the last calibration;
	
	\item[\textbf{-}] Reading uncertainty (particularly significant in the case of instruments with an analog output);
	
	\item[\textbf{-}] Instrumental load, i.e., alteration of the measured system resulting from interaction with the instrument;

	\item[\textbf{-}] Intrinsic uncertainty of the measurand, which is linked to its definition;
	
	\item[\textbf{-}] Uncertainty in the realization of the measurand’s definition, due, for example, to imperfect knowledge of the quantities that define the state of the measured system;

	\item[\textbf{-}] Uncertainty with which the influencing quantities are known.
\end{enumerate}

Another category of direct measurement methods involves a comparison by contrast between the quantity being measured and another quantity of the same type, which is generated by a variable standard. 

In this case, the measurement process involves an auxiliary device, whose function is to detect the condition of equivalence between the measurand $m$ and the quantity $c_1$ provided by the standard, so that the measurement is assigned as:

\begin{equation}
	m = c_1
\end{equation}

The uncertainty of the measured values obtained by applying a comparison method by opposition depends on the following factors:

\begin{enumerate}
	\item[\textbf{-}] Uncertainty and resolution of the sample;
	\item[\textbf{-}] The uncertainty with which the equivalence condition between the measurand and the quantity realized by the sample is satisfied;
	\item[\textbf{-}] Intrinsic uncertainty of the measurand;
	\item[\textbf{-}] Imperfect realization of the definition of the measurand;
	\item[\textbf{-}] Measurement uncertainty of influencing quantities.
\end{enumerate}

In electrical circuits, among the comparison methods, zero-crossing methods are often used, in which an auxiliary device is employed to detect a state of equilibrium, corresponding to the cancellation of a current or voltage in the circuit.  When the equilibrium condition is reached, the quantity under test $m$ can be expressed in terms of a number $n$ of known quantities provided by variable standards:

\begin{equation}
	m = f(c_1,c_2,\dots,c_n)
\end{equation}

In this case, the uncertainty contributions to be considered are:

\begin{enumerate}
	\item[\textbf{-}] Uncertainty and resolution of the samples;
	\item[\textbf{-}] Uncertainty associated with the realization of the zero condition;
	\item[\textbf{-}] Intrinsic uncertainty of the measurand;
	\item[\textbf{-}] Imperfect realization of the definition of the measurand;
	\item[\textbf{-}] Measurement uncertainty of influencing quantities.
\end{enumerate}


A variant of the zero method is the substitution method, which consists of two distinct phases.  

The first phase involves establishing the condition that allows the measurand to be expressed using the following relationship, employing devices characterized by high stability; as will be seen, their uncertainty does not affect the estimation of the measurand. 

In a second step, a variable sample that realizes a quantity $c_R$ homogeneous with the measurand is substituted into the measurement system; then, the equilibrium condition is reestablished by acting exclusively on the sample $c_R$ and leaving the other devices unchanged. Under these conditions, the following relationship holds:

\begin{equation}
	c_R = f(c_1,c_2,\dots,c_n)
\end{equation}

whereby the measurement $c_R$ is simply assigned as:

\begin{equation}
	m = c_R
\end{equation}

Note that, in this case, the measurement uncertainty $m$ depends on the uncertainty of the $c_R$ sample and on the short-term stability of the devices that generate the values $c_1,\;c_2,\;\dots, c_n\;$ which may be known only approximately.

\subsubsection{Indirect measurement methods}

For the purposes of this text, a measurement method is considered indirect when the value of a system’s parameter is determined through a calculation involving other parameters of the system itself or of other systems with which it interacts. 

These parameters are estimated using direct measurement methods. The application of an indirect measurement method therefore presupposes the existence of a mathematical model that explicitly expresses the relationship between the measurand and the other quantities subjected to direct measurement. 

If we denote the indirectly obtained measurand by $Y$ and the directly measured quantities by $X$, this relationship can be expressed by the equation:\\

\begin{equation}
	Y = f(X_1,X_2,\dots, X_N)
\end{equation}

where, for simplicity, the dependence on the quantities that define the system’s state in terms of measurement and on the influencing quantities has been omitted, although this dependence must always be taken into account.

Examples of indirect measurement include determining the volume of a spherical solid based on the measurement of its diameter, or estimating the electrical resistance of a device using the voltamperometric method.\\

In the case of indirect measurement methods, the uncertainty is estimated based on the uncertainties of the quantities obtained directly by applying the rules described in the following paragraphs. 

In this case, an additional contribution to uncertainty is provided by the so called model uncertainty, which takes into account the fact that the model expressed by the relationship just seen above does not adequately describe the interactions between the measurand and the other parameters subject to direct measurement. 

If we consider again the case of measuring the volume of a solid based on the measured diameter, the model uncertainty could be related, for example, to the fact that the solid is not exactly spherical or to imperfect knowledge of the thermal expansion law of the solid in question.

\subsection{Deterministic model}

The deterministic model for uncertainty estimation, which was widely used in the past and is still employed today in contexts where an overestimation of uncertainty is acceptable, involves assigning the measured value to a parameter in the form of a bounded interval, known as a confidence interval, which is typically symmetric about the mean value $m_0$ assigned to the measured parameter. 

The value range is characterized by the following properties:

\begin{enumerate}
	\item[\textbf{-}] It is reasonably certain that the measured quantity falls within the range of values;
	\item[\textbf{-}] All elements within the range of values are equally valid for representing the measured quantity.
\end{enumerate}

Even when the range of values is assigned by applying statistical techniques to a dataset obtained using a repeated-measurement method, all elements of the range of values must be considered equally representative of the measured quantity, since no assumptions are made regarding the probability distribution of the assigned values.\\

The measure m of a parameter can therefore be expressed using the following equation:\\

\begin{equation}
	m = (m_0 \pm \delta m)\; \unit{U} 
\end{equation}

where $\delta m$ is the absolute uncertainty of the measurement, expressed in the same unit of measurement U as $m_0$. The interval of width $2\delta m$ represents the range of values assigned as the measurement of the parameter.
When the ranges of values assigned using different methods or equipment, or at different times, to the same measurand and in the same system state have at least one element in common, then the measurements are said to be compatible. 
If one wishes to evaluate the compatibility between measurements referring to different states of the system, it is necessary to refer the measurements to a single reference state: in this case, compatibility also includes  the validity of the model describing the system under measurement.

In many cases, it is more meaningful to report the relative value of the uncertainty, or the relative percentage value, rather than the absolute value, using the following definitions:\\

\begin{equation}
	\varepsilon_m = \frac{\delta m}{m_0} \hspace{1cm} \varepsilon_{m\%} = \frac{\delta m}{m_0} \cdot 100
\end{equation}

Relative uncertainty, in fact, provides an immediate indication of the quality of a measurement, which may not be immediately apparent when reporting the absolute uncertainty. 

Consider, for example, two length measurements characterized by the same absolute uncertainty $\delta m$ of 1 mm but with two different measured values of 1 m and 100 m. In this case, using absolute uncertainty, the results of the two measurements would be reported as:\\



\begin{equation}
	m_1 = (1\; \pm\; 0.001) \; \unit{m} \hspace{1cm} m_2 = (100\; \pm\; 0.001) \; \unit{m}
\end{equation}

If, on the other hand, the measurements are given with the relative uncertainty expressed as a percentage, the result will be:\\

\begin{equation}
	m_1 = 1\; \unit{m} \pm\; 0.1 \%  \hspace{1cm} m_2 = 100\; \pm\; 0.001 \%
\end{equation}


This highlights the superior quality of the measurement $m_2$ compared to $m_2$.
In some cases, it is convenient to express the uncertainty as a reduced value, that is, relative to a conventional value mc of the measurand, using the following equation:\\

\begin{equation}
	\varepsilon_{red} = \frac{\delta m}{m_C}
\end{equation}

\subsubsection{Uncertainty propagation}

The measurement uncertainty of a parameter usually depends on multiple factors, as highlighted in the preceding paragraphs, and is obtained by appropriately combining the various uncertainty contributions.

In the case of measurements obtained using direct methods, one simply sums the various uncertainty contributions: in the case of a direct-reading measurement, for example, the instrumental uncertainty and the reading uncertainty will be summed, while in the case of a measurement assigned by applying a comparison method, the uncertainty associated with the known the sample used and the uncertainty associated with the imperfect realization of the condition of  equivalence between the measurand and the quantity generated by the sample.

If, on the other hand, the measurement is assigned using an indirect method, uncertainty contributions associated with the various quantities are combined according to the rules indicated below.\\

Let Y denote the measured parameter, which is related to other parameters $X_i$ by the following model:

\begin{equation}
	Y = f(X_1,X_2,\dots, X_N)
\end{equation}

The central value $y_0$ of the range assigned as a measure of $Y$ is obtained from the central values $x_{i0}$ of the other parameters as follows:

\begin{equation}
	y_0 = f(x_{10}, x_{20}, \dots, x_{N0})
\end{equation}

The absolute value of the measurement uncertainty, i.e., the half-width of the range of values for $y$, is instead obtained separately from the absolute uncertainties $\delta_{xi}$ of the quantities $x_i$ using the following equation:\\


\begin{equation}
	\delta_y = \bigg|  \frac{\partial f}{\partial x_1}   \bigg|_{(x_{10},\dots,x_{N0})} \cdot \delta_{x1} \; + \bigg| \frac{\partial f}{\partial x_N}   \bigg|_{(x_{10},\dots,x_{N0})} \cdot \delta_{xN} 
\end{equation}

In other terms we can write that:\\

\begin{equation}
	\delta_y = \sum_{i = 1}^N  \Bigg| \frac{\partial f}{\partial x_i} \bigg|_{x_0} \; \Bigg| \cdot \delta_{xi}	
\end{equation}

Where the therm $x_0$ in the previous formula stands for all of the $(x_{10},\dots,x_{N0})$. This is intended as the calculation of the $\delta_y$ as the sum of the absolute values of the partial derivative of the equation given calculated for the specific measure.\\

The model expressed by the first relation represents the approximate Taylor series expansion to the first-order terms (linear expansion) of the function $f$ around the point $(x_0, x_1, \dots, x_n)$. If the function is linear, the relation is exact, since terms of order higher than the first are all zero. \\

If, on the other hand, $f$ is nonlinear, the relation provides an approximate value of $\delta_y$; this approximation is usually acceptable, since the increments $\delta_{x1}$ to $\delta_{xN}$ (the uncertainties of the quantities $x_1$) are such that the contributions associated with terms of order higher than the first are negligible.\\

Note that the uncertainty propagation model expressed by the relation inevitably leads to an overestimation of the uncertainty of the quantity $y$, since it considers the maximum value of the uncertainties $\delta_{xi}$ and the various contributions are summed in absolute value; thus, the possibility of any compensation, even partial, between the uncertainty contributions is excluded a priori.

\subsubsection{Examples}

Below are some examples of uncertainty calculations based on the deterministic model for specific cases of the function $f$.

\subsubsection*{Sum and difference}

\begin{equation}
	Y = X_1 + X_2 - X_3
\end{equation}

Applying the previous equation, we obtain that:\\

\begin{equation}
	\delta_y = \delta_1 + \delta_2 - \delta_3
\end{equation}

From this, the following rule can be stated: \textbf{the absolute uncertainty of a quantity obtained as the sum and/or difference of other quantities $x_i$ is equal to the sum of the absolute uncertainties of the quantities $x_i$.}\\

In this case, the pessimistic nature of the deterministic model is clearly highlighted; according to this model, the individual contributions of uncertainty are such that they result in an overall uncertainty corresponding to the worst-case combination of their signs. 

In reality, however, it is possible that the nominal values $x_{10}$, $x_{20}$, and $x_{30}$ assigned to the quantities in the equation just seen may deviate from the actual values in such a way as to partially offset the uncertainty. This can occur, for example, if the deviations of $x_2$ and $x_3$ have the same sign, or if the deviation of $x_1$ has a sign opposite to that of $x_2$.\\

In the case of measurements obtained by difference, further analysis is warranted to guard against “unpleasant surprises.” 
Consider the quantity Y obtained as:\\

\begin{equation}
	Y = X_1 - X_2
\end{equation}

In this case the relative uncertainty is obtained as:\\

\begin{equation}
	\varepsilon_y = \frac{\delta_y}{y_0} = \frac{\delta_{x_1} + \delta_{x_2}}{x_{10} - x_{20}}
\end{equation}

If the values $x_{10}$ and $x_{20}$ are similar, the denominator of the expression just seen may become of the same order of magnitude or even smaller than the numerator, so that the relative uncertainty of the quantity obtained by difference can take on very high values. Consider, for example, the measurement of power losses $P$ in a 10 kW voltage transformer with an efficiency $\eta$ of 0.98, obtained as:

\begin{equation}
	P_p = P_i - P_u	
\end{equation}


where $P_i$ is the power absorbed by the primary winding of the transformer and $P_u$ is the power delivered to the secondary winding, both of which are measured with a relative uncertainty of 1\%. By applying the above equation, it is easy to verify that the relative uncertainty of the $P_p$ measurement is 99\%.

\subsubsection*{Product and quotient}

Let's now switch to a different case scenario:

\begin{equation}
	Y = \frac{X_1 \cdot X_2}{X_3} 
\end{equation}

The absolute uncertainty of the quantity $Y$ is given by:\\

\begin{equation}
	\delta_y = \frac{x_{20}}{x_{30}} \cdot \delta_{x_1} + \frac{x_{10}}{x_{30}} \cdot \delta_{x_2} + \frac{x_{10} \cdot x_{20}}{{x_{30}}^2} \cdot \delta_{x_3}
\end{equation}

that can be rewritten as:\\

\begin{equation}
	\delta_y = \frac{x_{10} \cdot x_{20}}{x_{30}} \cdot \frac{\delta_{x_1}}{x_{10}} + \frac{x_{10} \cdot x_{20}}{x_{30}} \cdot \frac{\delta_{x_2}}{x_{20}} + \frac{x_{10} \cdot x_{20}}{x_{30}} \cdot \frac{\delta_{x_3}}{x_{30}} = y_0 \cdot (\varepsilon_{x_1} + \varepsilon_{x_2} + \varepsilon_{x_3})
\end{equation}\\

from which we can extrapolate the relative uncertainty of $y$ as:

\begin{equation}
	\varepsilon_y = \frac{\delta_y}{y_0} = \varepsilon_{x_1} + \varepsilon_{x_2} + \varepsilon_{x_3}
\end{equation}\\

The following rule therefore applies: \textbf{the relative uncertainty of a quantity obtained as a product and/or quotient of other quantities $x_i$ is equal to the sum of the relative uncertainties of the quantities $x_i$.}

\subsubsection*{Power and root}

Let's now switch to a different case scenario considering the power and root operations.

\begin{equation}
	Y_1 = X_1^N = \prod_N X_1 	
\end{equation}

Applying the rule of products, we simply get:

\begin{equation}
	\varepsilon_{y_1} = N \cdot \varepsilon_{x_1}
\end{equation}


Furthermore, since the root can be written as a power, such as:

\begin{equation}
	Y_2 = \sqrt[M]{X_2} = {X_2}^{\frac{1}{M}}
\end{equation}

we then get that:

\begin{equation}
	\varepsilon_{y_2} = \frac{1}{M} \cdot \varepsilon_{x_2} 
\end{equation}

\subsubsection*{Small increment}

Let us now consider the quantity $Y$ obtained as:

\begin{equation}
	Y = X + \Delta_X \hspace{1cm} \unit{supposing\;that} \hspace{1cm} \Delta_X \ll X 
\end{equation}

and suppose that the quantity being measured is the quantity $Z$ defined as follows:

\begin{equation}
	Z = \frac{Y}{X} = 1 + \frac{\Delta_X}{X}	
\end{equation}

From the last equation and using the definition of relative uncertainty, we have that:

\begin{equation}
	\varepsilon_z = \varepsilon_{(1 + \Delta_x/x)} = \frac{\delta_{(1 + \Delta_x/x)}}{1 + \Delta_x/x} \approx \delta_{(1 + \Delta_x/x)} 
\end{equation}

where the approximation used stems from the fact that, by assumption, $\frac{\Delta_x}{x}  \ll 1$. 

Applying the rules described above yields the following relationship:

\begin{equation}
	\delta_{(1+\Delta_x / x)} = \delta_1 + \delta_{\frac{\Delta_x}{x}} = \delta_{\frac{\Delta_x}{x}} = \frac{\Delta_x}{x} \cdot \varepsilon_{\frac{\Delta_x}{x}} = \frac{\Delta_x}{x} \cdot (\varepsilon_{\Delta_x} + \varepsilon_x) 
\end{equation}

which ultimately allows us to express the relative uncertainty of the measurand as:

\begin{equation}
	\varepsilon_z = \frac{\Delta_x}{x} \cdot (\varepsilon_{\Delta_x} + \varepsilon_x) 
\end{equation}

Note, therefore, that since the term $\Delta_X / X$ is much smaller than unity, the quantity measured by $Z$ is obtained with a relative uncertainty that is much smaller than the sum of the relative uncertainties of the input quantities.

\subsection{Probabilistic model}

The current reference framework for measurement uncertainty evaluation is provided by European Experimental Standard ENV 13005 issued in 1999which constitutes the official standard issued by the \textit{"Comité International des Poids et Mesures"} (CIPM).\\

The primary objective of this standard is the establishment of a unified and internationally accepted framework capable of yielding a more realistic evaluation of measurement uncertainty compared to the deterministic approach. The latter typically results in an overestimation of uncertainty due to the adoption of overly conservative assumptions, as discussed in the preceding section.\\

The framework defined in ENV 13005 is grounded in a probabilistic formulation, wherein the measurand is modelled as a random variable ($r.v.$) endowed with an associated probability density function ($p.d.f.$). 

Within this formulation, the assigned value of the measurand is interpreted as an estimate of the expected value of its probability distribution. 

Under the assumption of unbiased estimation, the dispersion of repeated observations around the mean value is characterized by a zero-mean random variable, whose standard deviation—derived from the associated $p.d.f.$ is adopted as the quantitative descriptor of measurement uncertainty, herein referred to as the standard uncertainty.\\

The assumption of unbiasedness of the measurand estimate implies the absence of residual systematic effects affecting the measurement result. Typical sources of systematic error include improper instrument zeroing and measurement system perturbations induced by the interaction between the measurand and the measuring instrument (e.g., loading effects).

The application of the probabilistic framework to uncertainty evaluation therefore requires that:

\begin{enumerate}[label = \textbf{-}]
	\item All statistically significant systematic effects are identified by the operator;
	\item Identified systematic effects are appropriately modeled and corrected within the measurement model.
\end{enumerate}

In contrast to the deterministic framework, in which the assigned range of values for the measurand defines an interval within which the true value is assumed to be certainly contained, the probabilistic framework interprets this range as an interval associated with a specified coverage probability. 

Such an interval is referred to as a confidence interval, while the associated probability is termed the confidence level.\\

Let $x$ denote the measurand and $u(x)$ its standard uncertainty. 
The confidence interval is defined in terms of the expanded uncertainty $U(x)$, obtained as the product of the standard uncertainty and a coverage factor $k$. Denoting by $x_0$ the estimate of the measurand, the confidence interval is expressed as:

\begin{equation}
	x_0 \pm k \cdot u(x)
\end{equation}

i.e., bounded by $x_0 - k \cdot u(x)$ and x $x_0 + k \cdot u(x)$ The confidence level can be rigorously determined only when the probability density function of the measurand is fully specified.\\

Within the ENV 13005 framework, the methods employed for the evaluation of uncertainty contributions are classified into two categories:

\begin{enumerate}[label=\textbf{-}]
	\item \textbf{Type A} evaluation of uncertainty, based on a statistical, frequency-based approach;
	\item \textbf{Type B} evaluation of uncertainty, performed using methods other than the frequency-based approach, for instance relying on a prior available information.
\end{enumerate}


Type A uncertainty evaluation primarily allows the quantification of uncertainty contributions of a random nature, such as noise superimposed on the measurement signal and the instability of the state variables of the measurement system. This approach requires the availability of a set of observations of the measurand, typically obtained through repeated measurements using an appropriate measurement procedure.\\

Type B uncertainty evaluation, on the other hand, accounts for uncertainty contributions of a systematic nature (not to be confused with systematic effects, which can usually be corrected). Such contributions include, for example, discrepancies between the actual and nominal calibration functions of a measurement device. \\

This type of evaluation is generally based on information provided by third parties, such as the measured value of a reference standard reported in a calibration certificate, or the uncertainty specifications declared by the instrument manufacturer. As will be shown later, Type B evaluation often requires the assumption of an a priori probability distribution for the considered uncertainty contributions.\\


The following sections will present the general rules for uncertainty evaluation according to the two defined categories, followed by a description of the uncertainty propagation model. Finally, methods for estimating the confidence interval (or equivalently the coverage factor) associated with a specified confidence level will be introduced.\\

\subsubsection{Category A uncertainty assessment}
\label{uncA}


Let a generic quantity $X$ be measured using a repeated-measurement method, and let $x_1, x_2, x_3, \dots, x_N$ denote the $N$ observations (the so-called sample) obtained under nominally identical conditions. It is further assumed that the observations of the quantity $X$ are mutually uncorrelated.\\

The population mean (or expected value) $\mu$ is estimated by means of the empirical mean  $\overline{x}$ of the available sample.


\begin{equation}
	\overline{x} = \frac{1}{N} \sum_{k=1}^N x_k
\label{eq:233}
\end{equation}

 It should be noted that the sample mean  $\overline{x}$ is itself a random variable whose probability density function is asymptotically Gaussian (see the Central Limit Theorem in). Moreover, if the observations $x_k$ are Gaussian-distributed, then the sample mean exhibits a Gaussian distribution even for finite sample sizes.\\



The population variance $\sigma^2$ is estimated through the corrected empirical variance $s^2(x)$ of the sample:

\begin{equation}
	s^2(x) = \frac{1}{N - 1} \sum_{k=1}^N(x_k - \overline{x})^2
\end{equation}

The positive square root $s(x)$ of the empirical variance, expressed in the same units as the quantity $X$, is referred to as the experimental standard deviation and represents the dispersion of the individual observations around the sample mean $\overline{x}$.\\

\autoref{figVV} illustrates, for example, two sets of voltage output measurements of a power supply. Both datasets are characterized by the same empirical mean $\overline{V}$, but by different values of the experimental standard deviation: it is observed that the second dataset, which exhibits a lower dispersion around the mean, is associated with a smaller standard deviation $s_2(V)$ than that of the first dataset $s_1(V)$.\\

The experimental standard deviation therefore represents the uncertainty contribution associated with random effects. 
Consequently, when the quantity $X$ is estimated via a single measurement, the corresponding standard uncertainty is taken equal to the experimental standard deviation, namely:

\begin{equation}
	u_A(x) = s(x)
\end{equation}


It should be noted that, even in the case of single-measurement methods, the experimental standard deviation should always be estimated, in order to avoid neglecting uncertainty contributions that may be significant.

When the quantity $X$ is instead estimated using a repeated-measurement method, the expected value of $X$ is represented by the sample mean $\overline{x}$.\\


\begin{figure}[h!]
\begin{center}

\begin{tikzpicture}[transform shape]
	\draw[line width=0.61234pt] (33, -7) -- (37, -7);
	\draw[line width=0.61234pt] (33, -11) -- (37, -11);
	\draw[line width=0.61234pt] (33, -7) -- (33, -11);
	\draw[line width=0.61234pt] (37, -7) -- (37, -11);
	\draw[line width=0.61234pt] (39, -7) -- (43, -7);
	\draw[line width=0.61234pt] (39, -11) -- (43, -11);
	\draw[line width=0.61234pt] (39, -7) -- (39, -11);
	\draw[line width=0.61234pt] (43, -7) -- (43, -11);
	\node[shape=rectangle, minimum width=0.965cm, minimum height=0.965cm] at (38, -9){} node[anchor=center, align=center, text width=0.577cm, inner sep=6pt] at (38, -9){\large $\bar{V}$};
	\node[shape=rectangle, minimum width=0.965cm, minimum height=0.965cm] at (35, -6.25){} node[anchor=center, align=center, text width=0.577cm, inner sep=6pt] at (35, -6.25){\large $V_1(V)$};
	\node[shape=rectangle, minimum width=0.965cm, minimum height=0.965cm] at (41, -6.25){} node[anchor=center, align=center, text width=0.577cm, inner sep=6pt] at (41, -6.25){\large $V_2(V)$};
	\node[shape=rectangle, minimum width=0.465cm, minimum height=0.465cm] at (32, -7){} node[anchor=center, align=center, text width=0.077cm, inner sep=6pt] at (32, -7){10.02};
	\node[shape=rectangle, minimum width=0.465cm, minimum height=0.465cm] at (32, -8){} node[anchor=center, align=center, text width=0.077cm, inner sep=6pt] at (32, -8){10.01\\};
	\node[shape=rectangle, minimum width=0.465cm, minimum height=0.465cm] at (32, -9){} node[anchor=center, align=center, text width=0.077cm, inner sep=6pt] at (32, -9){10.00\\};
	\node[shape=rectangle, minimum width=0.465cm, minimum height=0.465cm] at (32, -10){} node[anchor=center, align=center, text width=0.077cm, inner sep=6pt] at (32, -10){9.99\\};
	\node[shape=rectangle, minimum width=0.465cm, minimum height=0.465cm] at (32, -11){} node[anchor=center, align=center, text width=0.077cm, inner sep=6pt] at (32, -11){9.98\\};
	\draw[line width=0.3pt] (33, -9) -- (37, -9);
	\draw[line width=0.3pt, dash pattern={on 0.3pt off 2.4pt}] (33, -8) -- (37, -8);
	\draw[line width=0.3pt, dash pattern={on 0.3pt off 2.4pt}] (33, -10) -- (37, -10);
	\draw[line width=0.3pt, dash pattern={on 0.3pt off 2.4pt}] (34, -7) -- (34, -11);
	\draw[line width=0.3pt, dash pattern={on 0.3pt off 2.4pt}] (35, -7) -- (35, -11);
	\draw[line width=0.3pt, dash pattern={on 0.3pt off 2.4pt}] (36, -7) -- (36, -11);
	\draw[line width=0.2pt] (39, -9) -- (43, -9);
	\draw[line width=0.3pt, dash pattern={on 0.3pt off 2.4pt}] (39, -8) -- (43, -8);
	\draw[line width=0.3pt, dash pattern={on 0.3pt off 2.4pt}] (39, -10) -- (43, -10);
	\draw[line width=0.3pt, dash pattern={on 0.3pt off 2.4pt}] (40, -7) -- (40, -11);
	\draw[line width=0.3pt, dash pattern={on 0.3pt off 2.4pt}] (41, -7) -- (41, -11);
	\draw[line width=0.3pt, dash pattern={on 0.3pt off 2.4pt}] (42, -7) -- (42, -11);
	\node[shape=rectangle, minimum width=0.465cm, minimum height=0.465cm] at (33, -11.75){} node[anchor=center, align=center, text width=0.077cm, inner sep=6pt] at (33, -11.75){0\\};
	\node[shape=rectangle, minimum width=0.465cm, minimum height=0.465cm] at (35, -11.75){} node[anchor=center, align=center, text width=0.077cm, inner sep=6pt] at (35, -11.75){10\\};
	\node[shape=rectangle, minimum width=0.465cm, minimum height=0.465cm] at (34, -11.75){} node[anchor=center, align=center, text width=0.077cm, inner sep=6pt] at (34, -11.75){5\\};
	\node[shape=rectangle, minimum width=0.465cm, minimum height=0.465cm] at (36, -11.75){} node[anchor=center, align=center, text width=0.077cm, inner sep=6pt] at (36, -11.75){15\\};
	\node[shape=rectangle, minimum width=0.465cm, minimum height=0.465cm] at (37, -11.75){} node[anchor=center, align=center, text width=0.077cm, inner sep=6pt] at (37, -11.75){20};
	\node[shape=rectangle, minimum width=0.465cm, minimum height=0.465cm] at (39, -11.75){} node[anchor=center, align=center, text width=0.077cm, inner sep=6pt] at (39, -11.75){0\\};
	\node[shape=rectangle, minimum width=0.465cm, minimum height=0.465cm] at (41, -11.75){} node[anchor=center, align=center, text width=0.077cm, inner sep=6pt] at (41, -11.75){10\\};
	\node[shape=rectangle, minimum width=0.465cm, minimum height=0.465cm] at (40, -11.75){} node[anchor=center, align=center, text width=0.077cm, inner sep=6pt] at (40, -11.75){5\\};
	\node[shape=rectangle, minimum width=0.465cm, minimum height=0.465cm] at (42, -11.75){} node[anchor=center, align=center, text width=0.077cm, inner sep=6pt] at (42, -11.75){15\\};
	\node[shape=rectangle, minimum width=0.465cm, minimum height=0.465cm] at (43, -11.75){} node[anchor=center, align=center, text width=0.077cm, inner sep=6pt] at (43, -11.75){20};
	\node[circ] at (33.25, -9.75){};
	\node[circ] at (33.5, -10.25){};
	\node[circ] at (33.75, -7.75){};
	\node[circ] at (34, -8.25){};
	\node[circ] at (34.25, -8.75){};
	\node[circ] at (35.5, -10){};
	\node[circ] at (34.489, -10.25){};
	\node[circ] at (35, -8.25){};
	\node[circ] at (34.779, -9.279){};
	\node[circ] at (34.779, -9.279){};
	\node[circ] at (35.25, -9.5){};
	\node[circ] at (35.75, -9.75){};
	\node[circ] at (36, -8.25){};
	\node[circ] at (36.25, -9.5){};
	\node[circ] at (36.5, -9.75){};
	\node[circ] at (36.75, -7.75){};
	\node[circ] at (39.5, -9.25){};
	\node[circ] at (41.75, -8.75){};
	\node[circ] at (39.25, -8.75){};
	\node[circ] at (40.25, -8.75){};
	\node[circ] at (42.25, -9.25){};
	\node[circ] at (41.25, -9.25){};
	\node[circ] at (40.5, -9.25){};
	\node[circ] at (39.75, -8.5){};
	\node[circ] at (41, -9){};
	\node[circ] at (41.5, -8.5){};
	\node[circ] at (42.029, -9.779){};
	\node[circ] at (42.5, -8.75){};
\end{tikzpicture}

\end{center}
\label{figVV}
\caption{Two series of repeated measurements of the output voltage of a power supply. The second series is characterized by a lower dispersion $s_2(V) < s_1(V)$}
\end{figure}

Consequently, it is necessary to provide an estimate of the dispersion associated with the mean value itself, which in general could be obtained by analyzing multiple estimates of the mean derived from different sets of repeated measurements. 

However, this approach is often impractical; therefore, the variance of the mean is estimated from the empirical variance of the sample as follows:

\begin{equation}
	s^2(\overline{x}) = \frac{s^2(x)}{N}
\end{equation}

The experimental standard deviation of the mean, i.e. the dispersion of different estimates of the sample mean $\overline{x}$ around the expected value, is then given by:

\begin{equation}
	s(\overline{x}) = \frac{s(x)}{\sqrt{N}}
\end{equation}

In the case of a quantity X estimated via repeated measurements, the standard uncertainty is therefore taken equal to the experimental standard deviation of the mean, namely:

\begin{equation}
	u(x) = s(\overline{x}) = \frac{s(x)}{\sqrt{N}}
\label{eq:237}
\end{equation}


In the case of a quantity $X$ estimated via repeated measurements, the standard uncertainty is therefore taken as equal to the experimental standard deviation of the mean, namely:

\begin{equation}
	u_A(x) = s(\overline{x})
\end{equation}

At this stage, a critical remark concerning relation (\autoref{eq:237}) is necessary, as it appears to suggest that the standard uncertainty may be arbitrarily reduced simply by increasing the number $N$ of observations of the measured random variable.\\

First, it should be noted that, as $N$ increases, the observation time increases proportionally. 
Consequently, in addition to higher measurement costs, it becomes increasingly difficult to ensure the stability of the measurand: there is a risk of tracking temporal variations of the measurand rather than performing repeated observations under nominally identical conditions. 

Under such circumstances, the empirical estimators previously defined lose their validity, as illustrated for instance in the left graph of \autoref{figXX}, where it is observed that the sample mean $\overline{V}$ is no longer a consistent estimator of the population mean.\\

Furthermore, as the measurement duration increases, the risk of correlation among observations also increases, for example due to periodic fluctuations of influence quantities, as shown in the example of right graph of \autoref{figXX}. Under these conditions, estimator (\autoref{eq:237}) is no longer applicable, since the observations $x_1, x_2, x_3, \dots, x_N$ are no longer statistically independent.\\


\begin{figure}[h!]
\begin{center}

\begin{tikzpicture}[transform shape]
	% Paths, nodes and wires:
	\draw[line width=0.61234pt] (33, -7) -- (38, -7);
	\draw[line width=0.61234pt] (33, -11) -- (38, -11);
	\draw[line width=0.61234pt] (33, -7) -- (33, -11);
	\draw[line width=0.61234pt] (38, -7) -- (38, -11);
	\node[shape=rectangle, minimum width=0.965cm, minimum height=0.965cm] at (35.5, -6.25){} node[anchor=center, align=center, text width=0.577cm, inner sep=6pt] at (35.5, -6.25){\large $V_1(V)$};
	\node[shape=rectangle, minimum width=0.965cm, minimum height=0.965cm] at (42.5, -6.25){} node[anchor=center, align=center, text width=0.577cm, inner sep=6pt] at (42.5, -6.25){\large $V_2(V)$};
	\node[shape=rectangle, minimum width=0.465cm, minimum height=0.465cm] at (32, -7){} node[anchor=center, align=center, text width=0.077cm, inner sep=6pt] at (32, -7){10.06\\};
	\node[shape=rectangle, minimum width=0.465cm, minimum height=0.465cm] at (32, -8){} node[anchor=center, align=center, text width=0.077cm, inner sep=6pt] at (32, -8){10.04\\};
	\node[shape=rectangle, minimum width=0.465cm, minimum height=0.465cm] at (32, -9){} node[anchor=center, align=center, text width=0.077cm, inner sep=6pt] at (32, -9){10.02\\};
	\node[shape=rectangle, minimum width=0.465cm, minimum height=0.465cm] at (32, -10){} node[anchor=center, align=center, text width=0.077cm, inner sep=6pt] at (32, -10){10.00\\};
	\node[shape=rectangle, minimum width=0.465cm, minimum height=0.465cm] at (32, -11){} node[anchor=center, align=center, text width=0.077cm, inner sep=6pt] at (32, -11){9.98\\};
	\draw[line width=0.3pt] (33, -9) -- (38, -9);
	\draw[line width=0.3pt, dash pattern={on 0.3pt off 2.4pt}] (33, -10) -- (38, -10);
	\draw[line width=0.3pt, dash pattern={on 0.3pt off 2.4pt}] (34, -7) -- (34, -11);
	\draw[line width=0.3pt, dash pattern={on 0.3pt off 2.4pt}] (36, -7) -- (36, -11);
	\node[shape=rectangle, minimum width=0.465cm, minimum height=0.465cm] at (33, -11.75){} node[anchor=center, align=center, text width=0.077cm, inner sep=6pt] at (33, -11.75){0\\};
	\node[shape=rectangle, minimum width=0.465cm, minimum height=0.465cm] at (35, -11.75){} node[anchor=center, align=center, text width=0.077cm, inner sep=6pt] at (35, -11.75){200\\};
	\node[shape=rectangle, minimum width=0.465cm, minimum height=0.465cm] at (34, -11.75){} node[anchor=center, align=center, text width=0.077cm, inner sep=6pt] at (34, -11.75){100};
	\node[shape=rectangle, minimum width=0.465cm, minimum height=0.465cm] at (36, -11.75){} node[anchor=center, align=center, text width=0.077cm, inner sep=6pt] at (36, -11.75){300\\};
	\node[shape=rectangle, minimum width=0.465cm, minimum height=0.465cm] at (37, -11.75){} node[anchor=center, align=center, text width=0.077cm, inner sep=6pt] at (37, -11.75){400};
	\draw[line width=0.3pt, dash pattern={on 0.3pt off 2.4pt}] (37, -7) -- (37, -11);
	\draw[line width=0.3pt, dash pattern={on 0.3pt off 2.4pt}] (35, -7) -- (35, -11);
	\draw[line width=0.61234pt] (40, -7) -- (45, -7);
	\draw[line width=0.61234pt] (40, -11) -- (45, -11);
	\draw[line width=0.61234pt] (40, -7) -- (40, -11);
	\draw[line width=0.61234pt] (45, -7) -- (45, -11);
	\draw[line width=0.3pt] (40, -9) -- (45, -9);
	\draw[line width=0.3pt, dash pattern={on 0.3pt off 2.4pt}] (40, -8) -- (45, -8);
	\draw[line width=0.3pt, dash pattern={on 0.3pt off 2.4pt}] (40, -10) -- (45, -10);
	\draw[line width=0.3pt, dash pattern={on 0.3pt off 2.4pt}] (41, -7) -- (41, -11);
	\draw[line width=0.3pt, dash pattern={on 0.3pt off 2.4pt}] (43, -7) -- (43, -11);
	\draw[line width=0.3pt, dash pattern={on 0.3pt off 2.4pt}] (44, -7) -- (44, -11);
	\draw[line width=0.3pt, dash pattern={on 0.3pt off 2.4pt}] (42, -7) -- (42, -11);
	\node[shape=rectangle, minimum width=0.465cm, minimum height=0.465cm] at (38, -11.75){} node[anchor=center, align=center, text width=0.077cm, inner sep=6pt] at (38, -11.75){500};
	\node[shape=rectangle, minimum width=0.465cm, minimum height=0.465cm] at (40, -11.75){} node[anchor=center, align=center, text width=0.077cm, inner sep=6pt] at (40, -11.75){0\\};
	\node[shape=rectangle, minimum width=0.465cm, minimum height=0.465cm] at (42, -11.75){} node[anchor=center, align=center, text width=0.077cm, inner sep=6pt] at (42, -11.75){200\\};
	\node[shape=rectangle, minimum width=0.465cm, minimum height=0.465cm] at (41, -11.75){} node[anchor=center, align=center, text width=0.077cm, inner sep=6pt] at (41, -11.75){100};
	\node[shape=rectangle, minimum width=0.465cm, minimum height=0.465cm] at (43, -11.75){} node[anchor=center, align=center, text width=0.077cm, inner sep=6pt] at (43, -11.75){300\\};
	\node[shape=rectangle, minimum width=0.465cm, minimum height=0.465cm] at (44, -11.75){} node[anchor=center, align=center, text width=0.077cm, inner sep=6pt] at (44, -11.75){400};
	\node[shape=rectangle, minimum width=0.465cm, minimum height=0.465cm] at (45, -11.75){} node[anchor=center, align=center, text width=0.077cm, inner sep=6pt] at (45, -11.75){500};
	\node[shape=rectangle, minimum width=0.465cm, minimum height=0.465cm] at (39, -7){} node[anchor=center, align=center, text width=0.077cm, inner sep=6pt] at (39, -7){10.02};
	\node[shape=rectangle, minimum width=0.465cm, minimum height=0.465cm] at (39, -8){} node[anchor=center, align=center, text width=0.077cm, inner sep=6pt] at (39, -8){10.01\\};
	\node[shape=rectangle, minimum width=0.465cm, minimum height=0.465cm] at (39, -9){} node[anchor=center, align=center, text width=0.077cm, inner sep=6pt] at (39, -9){10.00\\};
	\node[shape=rectangle, minimum width=0.465cm, minimum height=0.465cm] at (39, -10){} node[anchor=center, align=center, text width=0.077cm, inner sep=6pt] at (39, -10){9.99\\};
	\node[shape=rectangle, minimum width=0.465cm, minimum height=0.465cm] at (39, -11){} node[anchor=center, align=center, text width=0.077cm, inner sep=6pt] at (39, -11){9.98\\};
	\node[circ] at (33.25, -10.5){};
	\node[circ] at (33.5, -10.25){};
	\node[circ] at (33.5, -10.5){};
	\node[circ] at (33.25, -10){};
	\node[circ] at (33.5, -9.75){};
	\node[circ] at (33.5, -10){};
	\node[circ] at (33.5, -10.25){};
	\node[circ] at (33.8, -10){};
	\node[circ] at (33.8, -10.3){};
	\node[circ] at (33.5, -9.75){};
	\node[circ] at (33.75, -9.5){};
	\node[circ] at (33.7, -9.7){};
	\node[circ] at (34.1, -10.2){};
	\node[circ] at (34, -10.5){};
	\node[circ] at (34, -9.7){};
	\node[circ] at (34, -10){};
	\node[circ] at (34.25, -9.75){};
	\node[circ] at (34.25, -10){};
	\node[circ] at (34.25, -9.25){};
	\node[circ] at (34.25, -9.5){};
	\node[circ] at (34.5, -9){};
	\node[circ] at (34, -9.25){};
	\node[circ] at (34.4, -9.5){};
	\node[circ] at (34.6, -9.7){};
	\node[circ] at (34.5, -9.5){};
	\node[circ] at (35, -9.25){};
	\node[circ] at (34.75, -9.5){};
	\node[circ] at (34.7, -8.9){};
	\node[circ] at (35, -8.75){};
	\node[circ] at (35, -9){};
	\node[circ] at (35, -9.25){};
	\node[circ] at (35.25, -9){};
	\node[circ] at (35.3, -9.2){};
	\node[circ] at (35, -8.7){};
	\node[circ] at (35.2, -8.5){};
	\node[circ] at (35.25, -8.75){};
	\node[circ] at (35.5, -9.25){};
	\node[circ] at (35.25, -9.5){};
	\node[circ] at (35.4, -8.8){};
	\node[circ] at (35.5, -9){};
	\node[circ] at (35.75, -8.75){};
	\node[circ] at (35.75, -9){};
	\node[circ] at (35.75, -8.25){};
	\node[circ] at (35.75, -8.5){};
	\node[circ] at (35.5, -8){};
	\node[circ] at (35.5, -8.3){};
	\node[circ] at (35.75, -8.25){};
	\node[circ] at (36, -8.75){};
	\node[circ] at (35, -9.5){};
	\node[circ] at (36, -8){};
	\node[circ] at (36, -8.25){};
	\node[circ] at (36.25, -8.25){};
	\node[circ] at (36.25, -8){};
	\node[circ] at (36.5, -8){};
	\node[circ] at (36.75, -8){};
	\node[circ] at (36.75, -7.75){};
	\node[circ] at (37, -7.75){};
	\node[circ] at (37, -8){};
	\node[circ] at (37, -8){};
	\node[circ] at (37.25, -7.75){};
	\node[circ] at (37.25, -8){};
	\node[circ] at (37.5, -8){};
	\node[circ] at (37.75, -8){};
	\node[circ] at (37.75, -7.5){};
	\node[circ] at (37.75, -7.75){};
	\node[circ] at (37.5, -7.25){};
	\node[circ] at (37.5, -7.5){};
	\node[circ] at (37.75, -7.5){};
	\node[circ] at (36.4, -8.7){};
	\node[circ] at (36.5, -8.5){};
	\node[circ] at (36.75, -8.5){};
	\node[circ] at (34.6, -9.7){};
	\node[circ] at (34.6, -9.2){};
	\node[circ] at (34.6, -9.45){};
	\node[circ] at (34.6, -9.2){};
	\node[circ] at (34.95, -9.65){};
	\node[circ] at (34.85, -9.45){};
	\node[circ] at (36.15, -8.35){};
	\node[circ] at (36.3, -8.4){};
	\node[circ] at (36.65, -8.35){};
	\node[circ] at (36.65, -7.85){};
	\node[circ] at (36.65, -8.1){};
	\node[circ] at (36.4, -7.6){};
	\node[circ] at (36.4, -7.9){};
	\node[circ] at (36.65, -7.85){};
	\node[circ] at (33.4, -10.4){};
	\node[circ] at (33.55, -10.45){};
	\node[circ] at (33.9, -10.4){};
	\node[circ] at (33.9, -9.9){};
	\node[circ] at (33.9, -10.15){};
	\node[circ] at (33.65, -9.65){};
	\node[circ] at (33.65, -9.95){};
	\node[circ] at (33.9, -9.9){};
	\node[circ] at (34.1, -9.8){};
	\node[circ] at (34.25, -9.85){};
	\node[circ] at (34.6, -9.8){};
	\node[circ] at (34.6, -9.3){};
	\node[circ] at (34.6, -9.55){};
	\node[circ] at (34.35, -9.05){};
	\node[circ] at (34.35, -9.35){};
	\node[circ] at (34.6, -9.3){};
	\node[circ] at (35.2, -8.85){};
	\node[circ] at (35.35, -8.9){};
	\node[circ] at (35.7, -8.85){};
	\node[circ] at (35.7, -8.35){};
	\node[circ] at (35.7, -8.6){};
	\node[circ] at (35.45, -8.1){};
	\node[circ] at (35.45, -8.4){};
	\node[circ] at (35.7, -8.35){};
	\node[circ] at (37.1, -8.1){};
	\node[circ] at (37.25, -8.15){};
	\node[circ] at (37.6, -8.1){};
	\node[circ] at (37.6, -7.6){};
	\node[circ] at (37.6, -7.85){};
	\node[circ] at (37.35, -7.35){};
	\node[circ] at (37.35, -7.65){};
	\node[circ] at (37.6, -7.6){};
	\node[circ] at (36.4, -8.4){};
	\node[circ] at (36.55, -8.45){};
	\node[circ] at (36.9, -8.4){};
	\node[circ] at (36.9, -7.9){};
	\node[circ] at (36.9, -8.15){};
	\node[circ] at (36.65, -7.65){};
	\node[circ] at (36.65, -7.95){};
	\node[circ] at (36.9, -7.9){};
	\draw[line width=0.3pt, dash pattern={on 0.3pt off 2.4pt}] (33, -8) -- (38, -8);
	\node[circ] at (40.15, -10.25){};
	\node[circ] at (40.3, -10.7){};
	\node[circ] at (40.4, -10.5){};
	\node[circ] at (40.65, -10.5){};
	\node[circ] at (40.2, -10.4){};
	\node[circ] at (40.55, -10.35){};
	\node[circ] at (40.3, -10.4){};
	\node[circ] at (40.45, -10.45){};
	\node[circ] at (40.35, -9.85){};
	\node[circ] at (40.5, -10.3){};
	\node[circ] at (40.6, -10.1){};
	\node[circ] at (40.85, -10.1){};
	\node[circ] at (40.4, -10){};
	\node[circ] at (40.75, -9.95){};
	\node[circ] at (40.5, -10){};
	\node[circ] at (40.65, -10.05){};
	\node[circ] at (42.5, -8.15){};
	\node[circ] at (42.4, -7.95){};
	\node[circ] at (42.65, -7.95){};
	\node[circ] at (42.5, -8.15){};
	\node[circ] at (42.5, -7.9){};
	\node[circ] at (42.5, -8.25){};
	\node[circ] at (42.5, -7.75){};
	\node[circ] at (42.5, -8){};
	\node[circ] at (42.25, -7.8){};
	\node[circ] at (42.5, -7.75){};
	\node[circ] at (42.9, -8.25){};
	\node[circ] at (42.8, -8.05){};
	\node[circ] at (43.05, -8.05){};
	\node[circ] at (42.9, -8.25){};
	\node[circ] at (42.9, -8){};
	\node[circ] at (42.9, -8.35){};
	\node[circ] at (42.9, -7.85){};
	\node[circ] at (42.9, -8.1){};
	\node[circ] at (42.65, -7.9){};
	\node[circ] at (42.9, -7.85){};
	\node[circ] at (42.1, -8.35){};
	\node[circ] at (42, -8.15){};
	\node[circ] at (42.25, -8.15){};
	\node[circ] at (42.1, -8.35){};
	\node[circ] at (42.1, -8.1){};
	\node[circ] at (42.1, -8.45){};
	\node[circ] at (42.1, -7.95){};
	\node[circ] at (42.1, -8.2){};
	\node[circ] at (41.85, -8){};
	\node[circ] at (42.1, -7.95){};
	\node[circ] at (44.35, -10.65){};
	\node[circ] at (44.45, -10.45){};
	\node[circ] at (44.7, -10.45){};
	\node[circ] at (44.25, -10.35){};
	\node[circ] at (44.6, -10.3){};
	\node[circ] at (44.35, -10.35){};
	\node[circ] at (44.5, -10.4){};
	\node[circ] at (44.55, -10.25){};
	\node[circ] at (44.65, -10.05){};
	\node[circ] at (44.05, -10.25){};
	\node[circ] at (44.15, -10.05){};
	\node[circ] at (44.4, -10.05){};
	\node[circ] at (43.95, -9.95){};
	\node[circ] at (44.3, -9.9){};
	\node[circ] at (44.05, -9.95){};
	\node[circ] at (44.2, -10){};
	\node[circ] at (44.25, -9.85){};
	\node[circ] at (44.35, -9.65){};
	\node[circ] at (41.45, -8.45){};
	\node[circ] at (41.15, -9.05){};
	\node[circ] at (41.65, -8.75){};
	\node[circ] at (41.25, -9.25){};
	\node[circ] at (41.5, -9.05){};
	\node[circ] at (41.65, -8.35){};
	\node[circ] at (41.65, -8.75){};
	\node[circ] at (41.35, -8.75){};
	\node[circ] at (41.35, -9){};
	\node[circ] at (41.35, -8.75){};
	\node[circ] at (41.65, -8.85){};
	\node[circ] at (41.35, -8.85){};
	\node[circ] at (41.35, -9.1){};
	\node[circ] at (41.5, -8.6){};
	\node[circ] at (41.1, -8.9){};
	\node[circ] at (41.35, -8.85){};
	\node[circ] at (42.15, -8.15){};
	\node[circ] at (41.5, -8.3){};
	\node[circ] at (42.15, -8.35){};
	\node[circ] at (41.75, -8.85){};
	\node[circ] at (42, -8.65){};
	\node[circ] at (41.75, -7.95){};
	\node[circ] at (42.15, -8.35){};
	\node[circ] at (41.85, -8.35){};
	\node[circ] at (41.85, -8.6){};
	\node[circ] at (41.85, -8.35){};
	\node[circ] at (41.95, -8.35){};
	\node[circ] at (41.65, -8.35){};
	\node[circ] at (41.85, -8.7){};
	\node[circ] at (42, -8.2){};
	\node[circ] at (41.65, -8.25){};
	\node[circ] at (41.65, -8.35){};
	\node[circ] at (40.95, -9.5){};
	\node[circ] at (40.65, -10.1){};
	\node[circ] at (41.15, -9.8){};
	\node[circ] at (40.75, -10.3){};
	\node[circ] at (41, -10.1){};
	\node[circ] at (41.15, -9.4){};
	\node[circ] at (41.15, -9.8){};
	\node[circ] at (40.85, -9.8){};
	\node[circ] at (40.85, -10.05){};
	\node[circ] at (40.85, -9.8){};
	\node[circ] at (41.15, -9.9){};
	\node[circ] at (40.85, -9.9){};
	\node[circ] at (40.85, -10.15){};
	\node[circ] at (41, -9.65){};
	\node[circ] at (40.6, -9.95){};
	\node[circ] at (40.85, -9.9){};
	\node[circ] at (41.65, -9.2){};
	\node[circ] at (41, -9.35){};
	\node[circ] at (41.65, -9.4){};
	\node[circ] at (41.25, -9.9){};
	\node[circ] at (41.5, -9.7){};
	\node[circ] at (41.25, -9){};
	\node[circ] at (41.65, -9.4){};
	\node[circ] at (41.35, -9.4){};
	\node[circ] at (41.35, -9.65){};
	\node[circ] at (41.35, -9.4){};
	\node[circ] at (41.45, -9.4){};
	\node[circ] at (41.15, -9.4){};
	\node[circ] at (41.35, -9.75){};
	\node[circ] at (41.5, -9.25){};
	\node[circ] at (41.15, -9.3){};
	\node[circ] at (41.15, -9.4){};
	\node[circ, rotate=90] at (43.929, -9.571){};
	\node[circ, rotate=90] at (44.148, -9.966){};
	\node[circ, rotate=90] at (44.229, -9.371){};
	\node[circ, rotate=90] at (44.348, -9.866){};
	\node[circ, rotate=90] at (44.298, -9.566){};
	\node[circ, rotate=90] at (43.829, -9.371){};
	\node[circ, rotate=90] at (44.229, -9.371){};
	\node[circ, rotate=90] at (43.848, -9.766){};
	\node[circ, rotate=90] at (44.098, -9.766){};
	\node[circ, rotate=90] at (43.848, -9.766){};
	\node[circ, rotate=90] at (44.329, -9.371){};
	\node[circ, rotate=90] at (43.948, -9.766){};
	\node[circ, rotate=90] at (44.198, -9.766){};
	\node[circ, rotate=90] at (44.079, -9.521){};
	\node[circ, rotate=90] at (43.998, -10.016){};
	\node[circ, rotate=90] at (43.948, -9.766){};
	\node[circ, rotate=90] at (43.629, -8.871){};
	\node[circ, rotate=90] at (43.779, -9.521){};
	\node[circ, rotate=90] at (43.829, -8.871){};
	\node[circ, rotate=90] at (44.329, -9.271){};
	\node[circ, rotate=90] at (44.129, -9.021){};
	\node[circ, rotate=90] at (43.429, -9.271){};
	\node[circ, rotate=90] at (43.829, -8.871){};
	\node[circ, rotate=90] at (43.829, -9.171){};
	\node[circ, rotate=90] at (44.079, -9.171){};
	\node[circ, rotate=90] at (43.829, -9.171){};
	\node[circ, rotate=90] at (43.829, -9.071){};
	\node[circ, rotate=90] at (43.829, -9.371){};
	\node[circ, rotate=90] at (44.2, -9.2){};
	\node[circ, rotate=90] at (43.679, -9.021){};
	\node[circ, rotate=90] at (43.729, -9.371){};
	\node[circ, rotate=90] at (43.829, -9.371){};
	\node[circ, rotate=90] at (43.298, -8.866){};
	\node[circ, rotate=90] at (43.5, -9.3){};
	\node[circ, rotate=90] at (43.2, -8.8){};
	\node[circ, rotate=90] at (43.7, -9.2){};
	\node[circ, rotate=90] at (43.5, -8.95){};
	\node[circ, rotate=90] at (43.198, -8.666){};
	\node[circ, rotate=90] at (43.2, -8.8){};
	\node[circ, rotate=90] at (43.2, -9.1){};
	\node[circ, rotate=90] at (43.45, -9.1){};
	\node[circ, rotate=90] at (43.2, -9.1){};
	\node[circ, rotate=90] at (43.3, -8.8){};
	\node[circ, rotate=90] at (43.3, -9.1){};
	\node[circ, rotate=90] at (43.55, -9.1){};
	\node[circ, rotate=90] at (43.448, -8.816){};
	\node[circ, rotate=90] at (43.35, -9.35){};
	\node[circ, rotate=90] at (43.3, -9.1){};
	\node[circ, rotate=90] at (42.998, -8.166){};
	\node[circ, rotate=90] at (43.148, -8.816){};
	\node[circ, rotate=90] at (43.198, -8.166){};
	\node[circ, rotate=90] at (43.3, -8.7){};
	\node[circ, rotate=90] at (43.498, -8.316){};
	\node[circ, rotate=90] at (42.798, -8.566){};
	\node[circ, rotate=90] at (43.198, -8.166){};
	\node[circ, rotate=90] at (43.198, -8.466){};
	\node[circ, rotate=90] at (43.448, -8.466){};
	\node[circ, rotate=90] at (43.198, -8.466){};
	\node[circ, rotate=90] at (43.198, -8.366){};
	\node[circ, rotate=90] at (43.198, -8.666){};
	\node[circ, rotate=90] at (43.171, -8.629){};
	\node[circ, rotate=90] at (43.048, -8.316){};
	\node[circ, rotate=90] at (43.098, -8.666){};
	\node[circ, rotate=90] at (43.198, -8.666){};
\end{tikzpicture}

\end{center}	
\caption{Example of measurements series taken with an increased number of samples}
\label{figXX}
\end{figure}


With increasing number of observations, it therefore becomes increasingly difficult to guarantee the validity conditions underlying the empirical estimators previously introduced. It should also be emphasised that the considerations presented so far refer exclusively to uncertainty contributions of random nature, and that increasing the number of observations does not reduce systematic contributions, the most relevant of which, in most practical cases, is the instrumental uncertainty, i.e. the deviation of the instrument calibration function from its nominal behaviour.\\

It is therefore reasonable to select a number of observations $N$ such that the standard deviation of the mean (\autoref{eq:237}) becomes negligible, or at least of the same order of magnitude, with respect to the other uncertainty contributions.\\

A further relevant aspect concerns the estimation of the probability density function ($p.d.f.$) of the sample mean $\overline{x}$, which is obtained through relation (\autoref{eq:233}). As previously recalled in connection with the Central Limit Theorem, $\overline{x}$ is a random variable whose $p.d.f.$ is asymptotically Gaussian. 

Nevertheless, the sample mean is often assumed to follow a Gaussian distribution even for small sample sizes, including the case $N = 1$. 
This approximation is generally acceptable, also because an exact determination of the $p.d.f.$ would require the application of complex analytical methods based on the convolution of the $p.d.f.$s of the population of individual observations.\\

Such analytical effort would only be justified if it led to a significantly more accurate evaluation of the confidence interval associated with a prescribed confidence level (see \autoref{uncA}). 

However, the estimator of the experimental standard deviation of the mean, expressed by (\autoref{eq:237}), is itself not exact but affected by uncertainty, which may assume non-negligible values.\\

For instance, considering a normally distributed random variable $X$, the relative uncertainty of the estimator of the experimental standard deviation of the mean can be expressed as:

\begin{equation}
	\frac{\sigma[s(\overline{x})]}{s(\overline{x})} \approx \frac{1}{\sqrt{2 \cdot (N - 1)}}
\end{equation}

Using this relation, relative uncertainties of approximately 24\% for $N=10$ and 16\% for $N=20$ are obtained. 

It is therefore generally unjustified to pursue a more refined modeling of the $p.d.f.$ of the random variable $X$, since the confidence interval is ultimately derived from a parameter that is typically estimated with an uncertainty on the order of 10–20\%.\\

A final remark concerns the assumption of absence of correlation among observations of the measurand, which is satisfied in most measurement applications. 
However, in specific cases—such as frequency standard analysis and other long-term measurement processes—the observations may be correlated, leading to biased estimators of both the mean and the experimental standard deviation of the mean. In such situations, specialized statistical methods for correlated data must be employed.

\subsubsection{Category B uncertainty assessment}


The evaluation of Type B uncertainty is not carried out through the statistical analysis of a set of observations of the measured quantity, but rather on the basis of information provided by third parties. 
Such a situation occurs, for example, in a comparison by opposition method (\ref{221}), where the measurement equation includes the value of a material standard that is typically stated in a calibration certificate.\\

Another example is measurements obtained using the direct-reading method, where one of the main uncertainty contributions is related to the imperfect knowledge of the calibration function of the instrument used, as stated by the manufacturer in the form of instrumental uncertainty.\\

As will be shown in the following section, the uncertainty of a measurement is obtained by combining uncertainties derived from both Type A and Type B methods; therefore, the latter must be expressed in the form of a standard uncertainty, clearly following rules different from those presented in the previous section.\\

A rule of general validity for the treatment of Type B uncertainties cannot be formulated, mainly because in many cases it is necessary to rely on prior information regarding the characteristics of instruments and standards, their history, or their behavior with respect to influence quantities. 

This highlights the fact that the evaluation of Type B uncertainty is not free from subjective judgment by the operator, which arises from their experience and sensitivity.\\


In the following, a list of cases commonly encountered in most measurement applications will be provided, starting from the best-defined situations and concluding with those in which the operator’s judgment can significantly influence the estimation of the uncertainty of the measured quantity.\\

\subsubsection*{Known expanded uncertainty and coverage factor}

The expanded measurement uncertainty $U(x)$ is stated, accompanied by the statement that the declared measurement uncertainty was obtained by multiplying the standard uncertainty $u(x)$ by the coverage factor $k$. 
In this case, the standard uncertainty is simply obtained by dividing the expanded uncertainty $U(x)$ by the coverage factor $k$.\\

This is clearly the optimal situation, which in the future should become the most common practice. At present, uncertainty is declared in this way in calibration certificates issued by primary metrology institutes and like the SIT calibration centers (Italian Calibration System).\\


\subsubsection*{Expanded uncertainty and known confidence level}


The expanded measurement uncertainty $U(x)$ and the confidence level are stated; the latter is typically 90\%, 95\%, or 99\%. 

Under these conditions, the standard uncertainty $u(x)$ can only be obtained by assuming that the \textit{"uncertainty"} random variable is characterized by a specific probability density function ($p.d.f.$).

For instance, if the declared confidence level is 95\% and a Gaussian distribution is assumed, the standard uncertainty is obtained by dividing the expanded uncertainty by the coverage factor 1.96.

This way of reporting uncertainty has been in use for approximately twenty years and is now adopted by almost all manufacturers of instruments and standards in the field of metrology.\\

\subsubsection*{Uncertainty according to the deterministic model}

Uncertainty, expressed through the class index of an instrument or by means of appropriate formulas, is declared with reference to a deterministic model. 
In this framework, the available information consists of an interval of width $2 \cdot \delta_x$ (where $\delta_x$ is the absolute uncertainty) within which the measurand is assumed to lie.\\

Uncertainty is expressed in these terms for older measurement instrumentation, such as electromechanical devices, and often also for electronic instrumentation, both analogue and digital, of a general-purpose type.\\

In order to obtain information in terms of standard uncertainty, it is necessary to assume a reasonable probability density function ($p.d.f.$) for the quantity under consideration.

% FIGURE REF
If all values within the interval can be considered equally valid representations of the measurand, it is natural to assume a uniform PDF for the measured random variable (see Figure 2.4). In this case, the corresponding standard uncertainty is given by:

\begin{equation}
	u(x) = \frac{\delta_x}{\sqrt{3}}
\label{eq:243}	
\end{equation}

If, instead, additional information suggests that values near the centre of the interval are more likely than those at the boundaries, a different \textit{p.d.f.} than the uniform one may be assumed. 

\begin{figure}[h!]
\begin{center}

\begin{tikzpicture}[transform shape]
	\draw[-latex] (33, -8) -- (38, -8);
	\draw[-latex] (35.5, -8.3) -| (35.5, -5);
	\draw[line width=0.8pt] (33.5, -8) -- (33.5, -6.5) -- (37.5, -6.5) -| (37.5, -8);
	\node[shape=rectangle, minimum width=0.465cm, minimum height=0.465cm] at (35, -4.5){} node[anchor=center, align=center, text width=0.077cm, inner sep=6pt] at (35, -4.5){$p(x)$};
	\node[shape=rectangle, minimum width=0.465cm, minimum height=0.465cm] at (35.5, -8.75){} node[anchor=center, align=center, text width=0.077cm, inner sep=6pt] at (35.5, -8.75){0\\};
	\node[shape=rectangle, minimum width=0.465cm, minimum height=0.465cm] at (37.5, -8.5){} node[anchor=center, align=center, text width=0.077cm, inner sep=6pt] at (37.5, -8.5){$\delta x$};
	\node[shape=rectangle, minimum width=0.465cm, minimum height=0.465cm] at (33.5, -8.5){} node[anchor=center, align=center, text width=0.077cm, inner sep=6pt] at (33.5, -8.5){$- \delta x$};
	\node[shape=rectangle, minimum width=0.465cm, minimum height=0.465cm] at (36, -6){} node[anchor=center, align=center, text width=0.077cm, inner sep=6pt] at (36, -6){$1/2\delta x$};
	\draw[-latex] (39, -8) -- (44, -8);
	\draw[-latex] (41.5, -8.3) -| (41.5, -5);
	\node[shape=rectangle, minimum width=0.465cm, minimum height=0.465cm] at (41, -4.5){} node[anchor=center, align=center, text width=0.077cm, inner sep=6pt] at (41, -4.5){$p(x)$};
	\node[shape=rectangle, minimum width=0.465cm, minimum height=0.465cm] at (41.5, -8.75){} node[anchor=center, align=center, text width=0.077cm, inner sep=6pt] at (41.5, -8.75){0\\};
	\node[shape=rectangle, minimum width=0.465cm, minimum height=0.465cm] at (43.5, -8.5){} node[anchor=center, align=center, text width=0.077cm, inner sep=6pt] at (43.5, -8.5){$\delta x$};
	\node[shape=rectangle, minimum width=0.465cm, minimum height=0.465cm] at (39.5, -8.5){} node[anchor=center, align=center, text width=0.077cm, inner sep=6pt] at (39.5, -8.5){$- \delta x$};
	\draw[line width=0.8pt] (39.5, -8) -- (41, -6.5) -- (42, -6.5) -- (43.5, -8);
	\node[shape=rectangle, minimum width=0.465cm, minimum height=0.465cm] at (40.5, -6){} node[anchor=center, align=center, text width=0.077cm, inner sep=6pt] at (40.5, -6){$- \beta \delta x$};
	\node[shape=rectangle, minimum width=0.465cm, minimum height=0.465cm] at (42, -6){} node[anchor=center, align=center, text width=0.077cm, inner sep=6pt] at (42, -6){$\beta \delta x$};
	\node[shape=rectangle, minimum width=0.465cm, minimum height=0.465cm] at (35.5, -9.75){} node[anchor=center, align=center, text width=0.077cm, inner sep=6pt] at (35.5, -9.75){$\textbf{(a)}$};
	\node[shape=rectangle, minimum width=0.465cm, minimum height=0.465cm] at (41.5, -9.75){} node[anchor=center, align=center, text width=0.077cm, inner sep=6pt] at (41.5, -9.75){$\textbf{(b)}$};
\end{tikzpicture}
	
\end{center}	
\label{FigZZ}
\caption{A random variable characterized by a uniform (a) and trapezoidal (b) probability distribution}
\end{figure}


An example is the trapezoidal probability density function (see Figure 2.4(b)), characterized by a base width equal to $2 \cdot \beta \cdot \delta_x$, where the parameter $\beta$ can take values between 0 and 1 (note that for $\beta = 1$ the distribution is uniform, whereas for $\beta = 0$ it becomes triangular).
In this case, the standard uncertainty is obtained from the half-width $\delta_x$ of the interval as:
  
\begin{equation}
  	u(x) = \delta_x \sqrt{\frac{1 - \beta^2}{6}}
\end{equation}

\subsubsection{Uncertainty propagation}
\label{uncprop243}

In the simple case of a single direct-reading measurement, the standard uncertainty is obtained by combining in quadrature the various uncertainty contributions, such as instrumental uncertainty and reading uncertainty, all of which are evaluated using Type B methods.

\begin{equation}
	u(x) = \sqrt{u_{1B}^2(x) + u_{2B}^2(x) + \dots}
\end{equation}

If the measurement is instead obtained by applying a measurement method based on repeated readings, or by using a mathematical relationship linking the measurand to other quantities, the rules described below are applied.\\

Let $Y$ denote the general measurand, which is related to $M$ quantities $X_i$ through the relationship:

\begin{equation}
	Y = f \left( X_1, \dots, X_H, X_{H+1} \dots, X_M \right)
\end{equation}

(where the first $H$ quantities are obtained by direct methods with repeated readings, while the remaining $M-H$ quantities are obtained by single-reading direct methods or are provided by third parties).

It should be recalled that, for a correct application of the probabilistic model of uncertainty evaluation, the measurement model must account for all significant systematic effects; therefore, \autoref{eq:243} must include any correction terms related, for example, to instrument loading or to the effect of influence quantities.

For each quantity $X_1, \dots, X_H$, a set of observations is available, so (\autoref{eq:233}) can be used to estimate the sample means $\overline{x}_1, \dots, \overline{H}$ of the available datasets. For the quantities $X_{H+1}, \dots, X_{M}$, only a single value is available, $x_{(H+1)0}, \dots, x_{M0}$ obtained, for example, from a single instrument reading or from data reported in a calibration certificate.\\

The estimate $y_0$ of the measurand is obtained as:

\begin{equation}
	y_0 = f \left(\overline{x_1}, \dots, \overline{x_H}, x_{(H+1)0}, \dots, x_{M0} \right)
\label{eq:244}	
\end{equation}


Regarding the evaluation of the measurement uncertainty, known as the \textbf{combined} standard uncertainty, denoted by $u_c(y)$, the following relationship holds:

\begin{equation}
	u_c(y) = \sqrt{\sum_{i=1}^M \left( \frac{\partial f}{\partial x_i} \right)^2 \cdot u^2(x_1) + 2 \cdot \sum_{j=1}^{M-1} \sum_{l=j+1}^M \frac{\partial f}{\partial x_j} \frac{\partial f}{\partial x_l} \cdot u(x_j, x_l) }
\label{eq:245}	
\end{equation}

where $u(x_i)$ is the standard uncertainty of the generic quantity $X_i$ appearing in \autoref{eq:244}, estimated by applying the appropriate uncertainty evaluation technique (Type A and/or Type B), while $u(x_j, x_l)$, known as the covariance of the pair of quantities $x_j$ and $x_l$ estimates of the two quantities.\\

The statistical dependence between two random variables is usually expressed by a dimensionless parameter called the correlation coefficient, which can take values between -1 and +1 and is obtained by normalizing the covariance with respect to the product of the standard deviations of the two variables:

\begin{equation}
	\rho(x_j, x_l) = \frac{u(x_j, x_l)}{u(x_j) \cdot u(x_i)}
\end{equation}

The relationship in \autoref{eq:245} highlights the fact that the estimation of uncertainty according to the probabilistic model cannot disregard the evaluation of the correlation between the estimates of the measured quantities.\\

However, there are no generally valid rules that allow the covariance, or equivalently the correlation coefficient, between two measured quantities to be calculated, since this evaluation depends on knowledge of the physical phenomena involved and on the measurement procedures, both of which rely on the operator’s experience and judgement.\\

It therefore appears that the probabilistic model for measurement uncertainty evaluation is suitable only for experts in the field of metrology, although there are situations in which simplified models can be used, as described below.\\

\subsubsection*{Absence of correlation}

When measurements of two quantities are obtained at different times and using different instruments, it is reasonable to assume the absence of correlation between the two measurements \footnote{Note that two variables may be independent in the sense of mathematical analysis, but their measurements may be dependent in the statistical sense—for example, because they are obtained from the same sample.}.

In this case, both the correlation coefficient and the covariance between the two measurements are equal to zero. Consequently, the expression for the combined standard uncertainty $u_c(Y)$ becomes:

\begin{equation}
	u_c(y) = \sqrt{\sum_{i = 1}^M \left(  \frac{\partial f}{\partial x_i}\right) \cdot u^2(x_i)}
\end{equation}


\subsubsection*{Strong correlation}

When measurements of two quantities are obtained using the same instrument, within the same range, and under the same conditions for the influence quantities, it is reasonable to assume a high degree of correlation between the two measurements.\\

If the instrumental uncertainty is the dominant uncertainty contribution, it is reasonable to assume $\rho(x_j, x_l) = +1$ (perfect positive linear correlation), and therefore \autoref{eq:245} becomes:

\begin{equation}
	u_c(y) = \sqrt{\left[ \sum_{i=1}^M \frac{\partial f}{\partial x_1} \cdot u(x_i)  \right]^2} = \sum_{i=i}^M \frac{\partial f}{\partial x_i} \cdot u(x_i)
\label{eq:248}	
\end{equation}

It should be noted that \autoref{eq:248} has a form similar to that used for uncertainty evaluation under the deterministic model; however, it presents two important differences: the absolute values of the partial derivatives do not appear, and the half-widths ($\delta_{x_i}$) of the value intervals are replaced by the standard uncertainties ($u(x_i)$) of the various quantities.\\


\subsubsection*{Analytical estimation of correlation}

In the case of two generic quantities $X_j$ and $X_l$ obtained by measurement methods based on repeated readings, i.e. when $N$ statistically independent observations are available for each quantity, it is possible to estimate not only the sample means $\overline{x_j}$ and $\overline{x_l}$ (\autoref{eq:233}) and the standard deviations of the mean $s(\overline{x_j}$ and $s(\overline{x_l}$ (\autoref{eq:237}), but also to compute an estimate of the covariance between the sample means using the following relationship:

\begin{equation}
	s(x_j, x_l) = \frac{1}{N \cdot (N-1)} \sum_{i=1}^N \left( x_{ji} - \overline{x_j} \right) \cdot \left( x_{li} - \overline{x_l} \right)
\label{eq:249}	
\end{equation}

Starting from (\autoref{eq:249}), it is also possible to obtain an estimate of the correlation coefficient between the sample means as follows:

\begin{equation}
	\rho (\overline{x_j},\overline{x_l}) = \frac{s(\overline{x_j},\overline{x_l})}{s(\overline{x_j}) \cdot s(\overline{x_l}}
\end{equation}


The left graph of \autoref{figure5} shows an example of two uncorrelated quantities, for which the correlation coefficient between the sample means is approximately zero, whereas the right graph of \autoref{figure5} shows an example of strongly correlated quantities, where the correlation coefficient between the sample means is approximately 0.94.\\


\begin{figure}[h!]
\begin{center}

\begin{tikzpicture}[transform shape]
	\node[circ] at (47.3, -11.8){};
	\node[circ] at (47.55, -11.55){};
	\node[circ] at (47.55, -11.8){};
	\node[circ] at (47.3, -11.3){};
	\node[circ] at (47.55, -11.05){};
	\node[circ] at (47.55, -11.3){};
	\node[circ] at (47.55, -11.55){};
	\node[circ] at (47.85, -11.3){};
	\node[circ] at (47.85, -11.6){};
	\node[circ] at (47.55, -11.05){};
	\node[circ] at (47.8, -10.8){};
	\node[circ] at (47.75, -11){};
	\node[circ] at (48.15, -11.5){};
	\node[circ] at (48.05, -11.8){};
	\node[circ] at (48.05, -11){};
	\node[circ] at (48.05, -11.3){};
	\node[circ] at (48.3, -11.05){};
	\node[circ] at (48.3, -11.3){};
	\node[circ] at (48.5, -10){};
	\node[circ] at (48.3, -10.8){};
	\node[circ] at (48.75, -9.75){};
	\node[circ] at (48.25, -10){};
	\node[circ] at (48.45, -10.8){};
	\node[circ] at (48.85, -10.45){};
	\node[circ] at (48.75, -10.25){};
	\node[circ] at (49.25, -10){};
	\node[circ] at (49, -10.25){};
	\node[circ] at (48.95, -9.65){};
	\node[circ] at (49, -9.25){};
	\node[circ] at (49.25, -9.75){};
	\node[circ] at (49.25, -10){};
	\node[circ] at (49.5, -9.75){};
	\node[circ] at (49.55, -9.95){};
	\node[circ] at (49, -9.2){};
	\node[circ] at (49.2, -9){};
	\node[circ] at (49.25, -9.25){};
	\node[circ] at (49.75, -10){};
	\node[circ] at (49.5, -10.25){};
	\node[circ] at (49.4, -9.3){};
	\node[circ] at (49.75, -9.75){};
	\node[circ] at (49.9, -9.2){};
	\node[circ] at (49.75, -9){};
	\node[circ] at (49.9, -8.7){};
	\node[circ] at (49.9, -8.95){};
	\node[circ] at (49.65, -8.45){};
	\node[circ] at (49.65, -8.75){};
	\node[circ] at (49.9, -8.7){};
	\node[circ] at (50.15, -9.2){};
	\node[circ] at (49.25, -10.25){};
	\node[circ] at (50.15, -8.45){};
	\node[circ] at (50.15, -8.7){};
	\node[circ] at (50.4, -8.7){};
	\node[circ] at (50.4, -8.45){};
	\node[circ] at (50.5, -8){};
	\node[circ] at (50.75, -8){};
	\node[circ] at (50.75, -7.75){};
	\node[circ] at (51, -7.75){};
	\node[circ] at (51, -8){};
	\node[circ] at (51, -8){};
	\node[circ] at (51.25, -7.75){};
	\node[circ] at (51.25, -8){};
	\node[circ] at (51.5, -8){};
	\node[circ] at (51.75, -8){};
	\node[circ] at (51.75, -7.5){};
	\node[circ] at (51.75, -7.75){};
	\node[circ] at (51.5, -7.25){};
	\node[circ] at (51.5, -7.5){};
	\node[circ] at (51.75, -7.5){};
	\node[circ] at (50.4, -8.7){};
	\node[circ] at (50.5, -8.5){};
	\node[circ] at (50.75, -8.5){};
	\node[circ] at (48.85, -10.45){};
	\node[circ] at (48.85, -9.95){};
	\node[circ] at (48.85, -10.2){};
	\node[circ] at (48.85, -9.95){};
	\node[circ] at (49.2, -10.4){};
	\node[circ] at (49.1, -10.2){};
	\node[circ] at (50.3, -8.8){};
	\node[circ] at (50.3, -8.4){};
	\node[circ] at (50.65, -8.35){};
	\node[circ] at (50.65, -7.85){};
	\node[circ] at (50.65, -8.1){};
	\node[circ] at (50.4, -7.6){};
	\node[circ] at (50.4, -7.9){};
	\node[circ] at (50.65, -7.85){};
	\node[circ] at (47.45, -11.7){};
	\node[circ] at (47.6, -11.75){};
	\node[circ] at (47.95, -11.7){};
	\node[circ] at (47.95, -11.2){};
	\node[circ] at (47.95, -11.45){};
	\node[circ] at (47.7, -10.95){};
	\node[circ] at (47.7, -11.25){};
	\node[circ] at (47.95, -11.2){};
	\node[circ] at (48.15, -11.1){};
	\node[circ] at (48.3, -11.15){};
	\node[circ] at (48.85, -10.55){};
	\node[circ] at (48.85, -10.05){};
	\node[circ] at (48.85, -10.3){};
	\node[circ] at (48.6, -9.8){};
	\node[circ] at (48.4, -10.65){};
	\node[circ] at (48.85, -10.05){};
	\node[circ] at (49.45, -9.6){};
	\node[circ] at (49.6, -9.65){};
	\node[circ] at (49.85, -9.3){};
	\node[circ] at (49.85, -8.8){};
	\node[circ] at (49.85, -9.05){};
	\node[circ] at (49.6, -8.55){};
	\node[circ] at (49.85, -9.6){};
	\node[circ] at (49.85, -8.8){};
	\node[circ] at (51.1, -8.1){};
	\node[circ] at (51.25, -8.15){};
	\node[circ] at (51.6, -8.1){};
	\node[circ] at (51.6, -7.6){};
	\node[circ] at (51.6, -7.85){};
	\node[circ] at (51.35, -7.35){};
	\node[circ] at (51.35, -7.65){};
	\node[circ] at (51.6, -7.6){};
	\node[circ] at (50.4, -8.4){};
	\node[circ] at (50.55, -8.45){};
	\node[circ] at (50.9, -8.4){};
	\node[circ] at (50.9, -7.9){};
	\node[circ] at (50.9, -8.15){};
	\node[circ] at (50.65, -7.65){};
	\node[circ] at (50.65, -7.95){};
	\node[circ] at (50.9, -7.9){};
	\draw[line width=0.61234pt] (40, -12) -- (45, -12);
	\draw[line width=0.61234pt] (40, -7) -- (40, -12);
	\draw[line width=0.61234pt] (45, -7) -- (45, -12);
	\draw[line width=0.3pt, dash pattern={on 0.3pt off 2.4pt}] (41, -7) -- (41, -12);
	\draw[line width=0.3pt, dash pattern={on 0.3pt off 2.4pt}] (43, -7) -- (43, -12);
	\draw[line width=0.3pt, dash pattern={on 0.3pt off 2.4pt}] (44, -7) -- (44, -12);
	\draw[line width=0.3pt, dash pattern={on 0.3pt off 2.4pt}] (42, -7) -- (42, -12);
	\draw[line width=0.3pt, dash pattern={on 0.3pt off 2.4pt}] (40, -8) -- (45, -8);
	\draw[line width=0.3pt, dash pattern={on 0.3pt off 2.4pt}] (40, -9) -- (45, -9);
	\draw[line width=0.3pt, dash pattern={on 0.3pt off 2.4pt}] (40, -10) -- (45, -10);
	\draw[line width=0.3pt, dash pattern={on 0.3pt off 2.4pt}] (40, -11) -- (45, -11);
	\draw[line width=0.61234pt] (40, -7) -- (45, -7);
	\draw[line width=0.61234pt] (47, -12) -- (52, -12);
	\draw[line width=0.61234pt] (47, -7) -- (47, -12);
	\draw[line width=0.61234pt] (52, -7) -- (52, -12);
	\draw[line width=0.3pt, dash pattern={on 0.3pt off 2.4pt}] (48, -7) -- (48, -12);
	\draw[line width=0.3pt, dash pattern={on 0.3pt off 2.4pt}] (50, -7) -- (50, -12);
	\draw[line width=0.3pt, dash pattern={on 0.3pt off 2.4pt}] (51, -7) -- (51, -12);
	\draw[line width=0.3pt, dash pattern={on 0.3pt off 2.4pt}] (49, -7) -- (49, -12);
	\draw[line width=0.3pt, dash pattern={on 0.3pt off 2.4pt}] (47, -8) -- (52, -8);
	\draw[line width=0.3pt, dash pattern={on 0.3pt off 2.4pt}] (47, -9) -- (52, -9);
	\draw[line width=0.3pt, dash pattern={on 0.3pt off 2.4pt}] (47, -10) -- (52, -10);
	\draw[line width=0.3pt, dash pattern={on 0.3pt off 2.4pt}] (47, -11) -- (52, -11);
	\draw[line width=0.61234pt] (47, -7) -- (52, -7);
	\node[shape=rectangle, minimum width=0.465cm, minimum height=0.465cm] at (39.5, -9.5){} node[anchor=center, align=center, text width=0.077cm, inner sep=6pt] at (39.5, -9.5){\large $x_j$};
	\node[shape=rectangle, minimum width=0.465cm, minimum height=0.465cm] at (46.5, -9.5){} node[anchor=center, align=center, text width=0.077cm, inner sep=6pt] at (46.5, -9.5){\large $x_j$};
	\node[shape=rectangle, minimum width=0.465cm, minimum height=0.465cm] at (42.5, -12.5){} node[anchor=center, align=center, text width=0.077cm, inner sep=6pt] at (42.5, -12.5){\large $x_l$};
	\node[shape=rectangle, minimum width=0.465cm, minimum height=0.465cm] at (49.5, -12.5){} node[anchor=center, align=center, text width=0.077cm, inner sep=6pt] at (49.5, -12.5){\large $x_l$};
	\node[circ] at (48, -10.55){};
	\node[circ] at (48.05, -10.75){};
	\node[circ] at (48.25, -10.8){};
	\node[circ] at (48.25, -10.55){};
	\node[circ] at (47.95, -10.4){};
	\node[circ] at (48.1, -10.45){};
	\node[circ] at (48.35, -10.4){};
	\node[circ] at (50.1, -8.35){};
	\node[circ] at (50.15, -8.55){};
	\node[circ] at (50.35, -8.6){};
	\node[circ] at (50.35, -8.35){};
	\node[circ] at (50.05, -8.2){};
	\node[circ] at (50.2, -8.25){};
	\node[circ] at (50.45, -8.2){};
	\node[circ] at (49.45, -9.25){};
	\node[circ] at (49.5, -9.45){};
	\node[circ] at (49.7, -9.5){};
	\node[circ] at (49.7, -9.25){};
	\node[circ] at (49.4, -9.1){};
	\node[circ] at (49.55, -9.15){};
	\node[circ] at (49.8, -9.1){};
	\node[circ] at (48.4, -10.35){};
	\node[circ] at (48.45, -10.55){};
	\node[circ] at (48.65, -10.6){};
	\node[circ] at (48.65, -10.35){};
	\node[circ] at (48.35, -10.2){};
	\node[circ] at (48.5, -10.25){};
	\node[circ] at (48.75, -10.2){};
	\node[circ] at (49, -9.65){};
	\node[circ] at (49.05, -9.85){};
	\node[circ] at (49.25, -9.9){};
	\node[circ] at (49.25, -9.65){};
	\node[circ] at (48.95, -9.5){};
	\node[circ] at (49.1, -9.55){};
	\node[circ] at (49.35, -9.5){};
	\node[circ] at (51.2, -7.85){};
	\node[circ] at (51.25, -8.05){};
	\node[circ] at (51.45, -8.1){};
	\node[circ] at (51.45, -7.85){};
	\node[circ] at (51.15, -7.7){};
	\node[circ] at (51.3, -7.75){};
	\node[circ] at (51.55, -7.7){};
	\node[circ] at (41.85, -9.65){};
	\node[circ] at (42, -9.65){};
	\node[circ] at (42.3, -9.1){};
	\node[circ] at (41.95, -9.5){};
	\node[circ] at (41.8, -9.65){};
	\node[circ] at (41.8, -9.4){};
	\node[circ] at (41.9, -9.25){};
	\node[circ] at (41.95, -9.2){};
	\node[circ] at (42, -9.4){};
	\node[circ] at (42.2, -9.45){};
	\node[circ] at (42.2, -9.2){};
	\node[circ] at (41.9, -9.05){};
	\node[circ] at (42.05, -9.1){};
	\node[circ] at (42.3, -9.05){};
	\node[circ] at (42.25, -9.7){};
	\node[circ] at (41.25, -9.65){};
	\node[circ] at (41.55, -9.1){};
	\node[circ] at (41.2, -9.5){};
	\node[circ] at (42.2, -9.7){};
	\node[circ] at (42.45, -9.2){};
	\node[circ] at (41.15, -9.25){};
	\node[circ] at (41.2, -9.2){};
	\node[circ] at (41.25, -9.4){};
	\node[circ] at (41.45, -9.45){};
	\node[circ] at (41.45, -9.2){};
	\node[circ] at (42.429, -9.579){};
	\node[circ] at (41.3, -9.1){};
	\node[circ] at (42.829, -9.579){};
	\node[circ] at (42.25, -9.7){};
	\node[circ] at (41.25, -9.65){};
	\node[circ] at (41.2, -9.5){};
	\node[circ] at (42.2, -9.7){};
	\node[circ] at (41.95, -8.95){};
	\node[circ] at (41.25, -9.4){};
	\node[circ] at (41.45, -9.45){};
	\node[circ] at (43.5, -9.1){};
	\node[circ] at (43.55, -9.05){};
	\node[circ] at (41.679, -9.579){};
	\node[circ] at (42.979, -9.679){};
	\node[circ] at (43.05, -9.45){};
	\node[circ] at (43.2, -9.45){};
	\node[circ] at (42.829, -9.629){};
	\node[circ] at (43.15, -9.3){};
	\node[circ] at (43, -9.45){};
	\node[circ] at (43, -9.2){};
	\node[circ] at (43.1, -9.05){};
	\node[circ] at (42.479, -9.729){};
	\node[circ] at (43.2, -9.2){};
	\node[circ] at (43.4, -9.25){};
	\node[circ] at (42.729, -9.729){};
	\node[circ] at (42.429, -9.579){};
	\node[circ] at (42.579, -9.629){};
	\node[circ] at (42.829, -9.579){};
	\node[circ] at (43.45, -9.5){};
	\node[circ] at (42.079, -9.629){};
	\node[circ] at (43.4, -9.5){};
	\node[circ] at (42.979, -9.729){};
	\node[circ] at (43.629, -9.379){};
	\node[circ] at (44.029, -9.379){};
	\node[circ] at (43.45, -9.5){};
	\node[circ] at (43.4, -9.5){};
	\node[circ] at (42.979, -9.729){};
	\node[circ] at (44.029, -9.629){};
	\node[circ] at (44.079, -9.579){};
	\node[circ] at (42.879, -9.379){};
	\node[circ] at (44.179, -9.479){};
	\node[circ] at (41.471, -8.771){};
	\node[circ] at (41.621, -8.771){};
	\node[circ] at (41.421, -8.771){};
	\node[circ] at (41.871, -8.821){};
	\node[circ] at (41.821, -8.821){};
	\node[circ] at (42.05, -8.7){};
	\node[circ] at (42.45, -8.7){};
	\node[circ] at (41.871, -8.821){};
	\node[circ] at (41.821, -8.821){};
	\node[circ] at (42.6, -8.8){};
	\node[circ] at (42.45, -8.75){};
	\node[circ] at (42.1, -8.85){};
	\node[circ] at (42.35, -8.85){};
	\node[circ] at (42.05, -8.7){};
	\node[circ] at (42.2, -8.75){};
	\node[circ] at (42.45, -8.7){};
	\node[circ] at (41.7, -8.75){};
	\node[circ] at (42.6, -8.85){};
	\node[circ] at (42.6, -8.85){};
	\node[circ] at (43.65, -8.75){};
	\node[circ] at (43.7, -8.7){};
	\node[circ] at (42.621, -10.271){};
	\node[circ] at (42.771, -10.271){};
	\node[circ] at (42.721, -10.121){};
	\node[circ] at (42.571, -10.271){};
	\node[circ] at (42.571, -10.021){};
	\node[circ] at (42.771, -10.021){};
	\node[circ] at (42.812, -9.663){};
	\node[circ] at (42.862, -9.912){};
	\node[circ] at (42.812, -9.913){};
	\node[circ] at (43.042, -9.792){};
	\node[circ] at (43.279, -9.829){};
	\node[circ] at (42.862, -9.912){};
	\node[circ] at (42.812, -9.913){};
	\node[circ] at (42.45, -10.2){};
	\node[circ] at (43.429, -9.929){};
	\node[circ] at (43.5, -9.7){};
	\node[circ] at (43.65, -9.7){};
	\node[circ] at (43.279, -9.879){};
	\node[circ] at (43.45, -9.7){};
	\node[circ] at (43.092, -9.942){};
	\node[circ] at (43.342, -9.942){};
	\node[circ] at (43.042, -9.792){};
	\node[circ] at (43.192, -9.842){};
	\node[circ] at (43.279, -9.829){};
	\node[circ] at (42.85, -10.25){};
	\node[circ] at (43.85, -9.75){};
	\node[circ] at (43.429, -9.979){};
	\node[circ] at (43.85, -9.75){};
	\node[circ] at (43.429, -9.979){};
	\node[circ] at (43.329, -9.629){};
	\node[circ] at (42.4, -9.35){};
	\node[circ] at (42.35, -9.35){};
	\node[circ] at (42.579, -9.229){};
	\node[circ] at (42.979, -9.229){};
	\node[circ] at (42.4, -9.35){};
	\node[circ] at (42.35, -9.35){};
	\node[circ] at (43.129, -9.329){};
	\node[circ] at (42.979, -9.279){};
	\node[circ] at (42.629, -9.379){};
	\node[circ] at (42.879, -9.379){};
	\node[circ] at (42.579, -9.229){};
	\node[circ] at (42.729, -9.279){};
	\node[circ] at (42.979, -9.229){};
	\node[circ] at (43.129, -9.379){};
	\node[circ] at (43.129, -9.379){};
	\node[circ] at (42.9, -8.45){};
	\node[circ] at (42.5, -8.6){};
	\node[circ] at (42.55, -8.55){};
	\node[circ] at (42.8, -8.55){};
	\node[circ] at (42.5, -8.4){};
	\node[circ] at (42.65, -8.45){};
	\node[circ] at (42.9, -8.4){};
	\node[circ] at (42.15, -8.45){};
	\node[circ] at (43.05, -8.55){};
	\node[circ] at (41.75, -8.6){};
	\node[circ] at (41.8, -8.55){};
	\node[circ] at (42.05, -8.55){};
	\node[circ] at (41.9, -8.45){};
	\node[circ] at (43.05, -8.55){};
	\node[circ] at (42.071, -8.121){};
	\node[circ] at (42.221, -8.121){};
	\node[circ] at (42.021, -8.121){};
	\node[circ] at (42.471, -8.171){};
	\node[circ] at (42.421, -8.171){};
	\node[circ] at (42.65, -8.05){};
	\node[circ] at (43.05, -8.05){};
	\node[circ] at (42.471, -8.171){};
	\node[circ] at (42.421, -8.171){};
	\node[circ] at (43.2, -8.15){};
	\node[circ] at (43.05, -8.1){};
	\node[circ] at (42.7, -8.2){};
	\node[circ] at (42.95, -8.2){};
	\node[circ] at (42.65, -8.05){};
	\node[circ] at (42.8, -8.1){};
	\node[circ] at (43.05, -8.05){};
	\node[circ] at (42.3, -8.1){};
	\node[circ] at (43.2, -8.2){};
	\node[circ] at (43.2, -8.2){};
	\node[circ] at (43, -8.7){};
	\node[circ] at (42.95, -8.7){};
	\node[circ] at (43.179, -8.579){};
	\node[circ] at (43, -8.7){};
	\node[circ] at (42.95, -8.7){};
	\node[circ] at (43.179, -8.579){};
	\node[circ] at (42.379, -10.179){};
	\node[circ] at (42.529, -10.179){};
	\node[circ] at (42.479, -10.029){};
	\node[circ] at (42.329, -10.179){};
	\node[circ] at (42.329, -9.929){};
	\node[circ] at (42.429, -9.779){};
	\node[circ] at (42.479, -9.729){};
	\node[circ] at (42.529, -9.929){};
	\node[circ] at (42.579, -9.629){};
	\node[circ] at (41.779, -10.179){};
	\node[circ] at (42.079, -9.629){};
	\node[circ] at (41.729, -10.029){};
	\node[circ] at (41.679, -9.779){};
	\node[circ] at (41.729, -9.729){};
	\node[circ] at (41.779, -9.929){};
	\node[circ] at (41.979, -9.979){};
	\node[circ] at (41.979, -9.729){};
	\node[circ] at (41.829, -9.629){};
	\node[circ] at (41.779, -10.179){};
	\node[circ] at (41.729, -10.029){};
	\node[circ] at (41.779, -9.929){};
	\node[circ] at (41.979, -9.979){};
	\node[circ] at (42.208, -10.108){};
	\node[circ] at (42.608, -10.158){};
	\node[circ] at (43.05, -8.95){};
	\node[circ] at (43.2, -8.95){};
	\node[circ] at (43.15, -8.8){};
	\node[circ] at (43, -8.95){};
	\node[circ] at (43, -8.7){};
	\node[circ] at (43.2, -8.7){};
	\node[circ] at (43.4, -8.75){};
	\node[circ] at (43.45, -9){};
	\node[circ] at (43.4, -9){};
	\node[circ] at (43.629, -8.879){};
	\node[circ] at (43.708, -8.508){};
	\node[circ] at (43.45, -9){};
	\node[circ] at (43.4, -9){};
	\node[circ] at (42.879, -8.879){};
	\node[circ] at (43.858, -8.608){};
	\node[circ] at (43.708, -8.558){};
	\node[circ] at (43.679, -9.029){};
	\node[circ] at (43.929, -9.029){};
	\node[circ] at (43.629, -8.879){};
	\node[circ] at (43.779, -8.929){};
	\node[circ] at (43.708, -8.508){};
	\node[circ] at (43.279, -8.929){};
	\node[circ] at (43.858, -8.658){};
	\node[circ] at (43.858, -8.658){};
	\node[circ] at (42.808, -8.858){};
	\node[circ] at (42.958, -8.858){};
	\node[circ] at (42.908, -8.708){};
	\node[circ] at (42.758, -8.858){};
	\node[circ] at (42.758, -8.608){};
	\node[circ] at (42.858, -8.458){};
	\node[circ] at (42.958, -8.608){};
	\node[circ] at (43.037, -8.837){};
	\node[circ] at (43.2, -9.7){};
	\node[circ] at (43.35, -9.8){};
	\node[circ] at (43.579, -9.979){};
	\node[circ] at (43.729, -9.979){};
	\node[circ] at (43.2, -9.75){};
	\node[circ] at (43.529, -9.979){};
	\node[circ] at (43.2, -9.7){};
	\node[circ] at (43.821, -9.621){};
	\node[circ] at (43.771, -9.621){};
	\node[circ] at (43.35, -9.85){};
	\node[circ] at (43.821, -9.621){};
	\node[circ] at (43.771, -9.621){};
	\node[circ] at (43.35, -9.85){};
	\node[circ] at (44.4, -9.75){};
	\node[circ] at (44.45, -9.7){};
	\node[circ] at (44.55, -9.6){};
	\node[circ] at (43.342, -10.192){};
	\node[circ] at (43.392, -10.442){};
	\node[circ] at (43.342, -10.442){};
	\node[circ] at (43.571, -10.321){};
	\node[circ] at (43.65, -9.95){};
	\node[circ] at (43.392, -10.442){};
	\node[circ] at (43.342, -10.442){};
	\node[circ] at (43.8, -10.05){};
	\node[circ] at (43.871, -9.821){};
	\node[circ] at (44.021, -9.821){};
	\node[circ] at (43.65, -10){};
	\node[circ] at (43.821, -9.821){};
	\node[circ] at (43.621, -10.471){};
	\node[circ] at (43.871, -10.471){};
	\node[circ] at (43.571, -10.321){};
	\node[circ] at (43.721, -10.371){};
	\node[circ] at (43.65, -9.95){};
	\node[circ] at (43.221, -10.371){};
	\node[circ] at (44.221, -9.871){};
	\node[circ] at (43.8, -10.1){};
	\node[circ] at (44.221, -9.871){};
	\node[circ] at (43.8, -10.1){};
	\node[circ] at (43.7, -9.75){};
	\node[circ] at (41.821, -10.321){};
	\node[circ] at (41.971, -10.321){};
	\node[circ] at (42.271, -9.771){};
	\node[circ] at (41.921, -10.171){};
	\node[circ] at (41.771, -10.321){};
	\node[circ] at (41.771, -10.071){};
	\node[circ] at (41.871, -9.921){};
	\node[circ] at (41.921, -9.871){};
	\node[circ] at (41.971, -10.071){};
	\node[circ] at (42.171, -10.121){};
	\node[circ] at (42.171, -9.871){};
	\node[circ] at (41.871, -9.721){};
	\node[circ] at (42.021, -9.771){};
	\node[circ] at (42.271, -9.721){};
	\node[circ] at (42.221, -10.371){};
	\node[circ] at (41.221, -10.321){};
	\node[circ] at (41.521, -9.771){};
	\node[circ] at (41.171, -10.171){};
	\node[circ] at (42.171, -10.371){};
	\node[circ] at (42.421, -9.871){};
	\node[circ] at (41.121, -9.921){};
	\node[circ] at (41.171, -9.871){};
	\node[circ] at (41.221, -10.071){};
	\node[circ] at (41.421, -10.121){};
	\node[circ] at (41.421, -9.871){};
	\node[circ] at (42.4, -10.25){};
	\node[circ] at (41.271, -9.771){};
	\node[circ] at (42.8, -10.25){};
	\node[circ] at (42.221, -10.371){};
	\node[circ] at (41.221, -10.321){};
	\node[circ] at (41.171, -10.171){};
	\node[circ] at (42.171, -10.371){};
	\node[circ] at (42.421, -9.871){};
	\node[circ] at (41.221, -10.071){};
	\node[circ] at (41.421, -10.121){};
	\node[circ] at (41.65, -10.25){};
	\node[circ] at (42.8, -10.3){};
	\node[circ] at (42.45, -10.4){};
	\node[circ] at (42.7, -10.4){};
	\node[circ] at (42.4, -10.25){};
	\node[circ] at (42.55, -10.3){};
	\node[circ] at (42.8, -10.25){};
	\node[circ] at (42.05, -10.3){};
	\node[circ] at (42.85, -10.05){};
	\node[circ] at (42.591, -10.942){};
	\node[circ] at (42.741, -10.942){};
	\node[circ] at (42.692, -10.792){};
	\node[circ] at (42.542, -10.942){};
	\node[circ] at (42.542, -10.692){};
	\node[circ] at (42.741, -10.692){};
	\node[circ] at (42.783, -10.333){};
	\node[circ] at (42.833, -10.583){};
	\node[circ] at (42.783, -10.583){};
	\node[circ] at (42.833, -10.583){};
	\node[circ] at (42.783, -10.583){};
	\node[circ] at (42.421, -10.871){};
	\node[circ] at (42.821, -10.921){};
	\node[circ] at (42.371, -10.021){};
	\node[circ] at (42.321, -10.021){};
	\node[circ] at (42.55, -9.9){};
	\node[circ] at (42.371, -10.021){};
	\node[circ] at (42.321, -10.021){};
	\node[circ] at (42.6, -10.05){};
	\node[circ] at (42.85, -10.05){};
	\node[circ] at (42.55, -9.9){};
	\node[circ] at (42.7, -9.95){};
	\node[circ] at (42.35, -10.85){};
	\node[circ] at (42.5, -10.85){};
	\node[circ] at (42.45, -10.7){};
	\node[circ] at (42.3, -10.85){};
	\node[circ] at (42.3, -10.6){};
	\node[circ] at (42.4, -10.45){};
	\node[circ] at (42.45, -10.4){};
	\node[circ] at (42.5, -10.6){};
	\node[circ] at (42.55, -10.3){};
	\node[circ] at (41.75, -10.85){};
	\node[circ] at (42.05, -10.3){};
	\node[circ] at (41.7, -10.7){};
	\node[circ] at (41.65, -10.45){};
	\node[circ] at (41.7, -10.4){};
	\node[circ] at (41.75, -10.6){};
	\node[circ] at (41.95, -10.65){};
	\node[circ] at (41.95, -10.4){};
	\node[circ] at (41.8, -10.3){};
	\node[circ] at (41.75, -10.85){};
	\node[circ] at (41.7, -10.7){};
	\node[circ] at (41.75, -10.6){};
	\node[circ] at (41.95, -10.65){};
	\node[circ] at (42.179, -10.779){};
	\node[circ] at (42.579, -10.829){};
	\node[circ] at (42.85, -9.55){};
	\node[circ] at (42.971, -10.571){};
	\node[circ] at (43.021, -10.821){};
	\node[circ] at (42.971, -10.821){};
	\node[circ] at (43.2, -10.7){};
	\node[circ] at (43.021, -10.821){};
	\node[circ] at (42.971, -10.821){};
	\node[circ] at (43.25, -10.85){};
	\node[circ] at (43.5, -10.85){};
	\node[circ] at (43.2, -10.7){};
	\node[circ] at (43.35, -10.75){};
	\node[circ] at (42.85, -10.75){};
	\node[circ] at (41.121, -8.671){};
	\node[circ] at (41.271, -8.671){};
	\node[circ] at (41.071, -8.671){};
	\node[circ] at (41.521, -8.721){};
	\node[circ] at (41.471, -8.721){};
	\node[circ] at (41.7, -8.6){};
	\node[circ] at (41.521, -8.721){};
	\node[circ] at (41.471, -8.721){};
	\node[circ] at (41.75, -8.75){};
	\node[circ] at (41.7, -8.6){};
	\node[circ] at (41.85, -8.65){};
	\node[circ] at (41.35, -8.65){};
	\node[circ] at (41.8, -8.35){};
	\node[circ] at (41.4, -8.5){};
	\node[circ] at (41.45, -8.45){};
	\node[circ] at (41.7, -8.45){};
	\node[circ] at (41.55, -8.35){};
	\node[circ] at (41.721, -8.021){};
	\node[circ] at (41.671, -8.021){};
	\node[circ] at (41.029, -9.329){};
	\node[circ] at (40.679, -9.429){};
	\node[circ] at (40.929, -9.429){};
	\node[circ] at (40.779, -9.329){};
	\node[circ] at (40.95, -9){};
	\node[circ] at (41.1, -9){};
	\node[circ] at (40.9, -9){};
	\node[circ] at (41.179, -8.979){};
	\node[circ] at (41.229, -8.829){};
	\node[circ] at (41.279, -8.779){};
	\node[circ] at (40.6, -8.9){};
	\node[circ] at (40.75, -8.9){};
	\node[circ] at (40.55, -8.9){};
	\node[circ] at (41, -8.95){};
	\node[circ] at (40.95, -8.95){};
	\node[circ] at (41.179, -8.829){};
	\node[circ] at (41, -8.95){};
	\node[circ] at (40.95, -8.95){};
	\node[circ] at (41.229, -8.979){};
	\node[circ] at (41.179, -8.829){};
	\node[circ] at (40.829, -8.879){};
	\node[circ] at (41.279, -8.579){};
	\node[circ] at (40.879, -8.729){};
	\node[circ] at (40.929, -8.679){};
	\node[circ] at (41.179, -8.679){};
	\node[circ] at (41.029, -8.579){};
	\node[circ] at (43.3, -10.7){};
	\node[circ] at (43.45, -10.7){};
	\node[circ] at (43.4, -10.55){};
	\node[circ] at (43.25, -10.7){};
	\node[circ] at (43.25, -10.45){};
	\node[circ] at (43.45, -10.45){};
	\node[circ] at (43.542, -10.342){};
	\node[circ] at (43.492, -10.342){};
	\node[circ] at (43.542, -10.342){};
	\node[circ] at (43.492, -10.342){};
	\node[circ] at (43.129, -10.629){};
	\node[circ] at (43.529, -10.679){};
	\node[circ] at (43.208, -10.608){};
	\node[circ] at (43.158, -10.458){};
	\node[circ] at (43.208, -10.358){};
	\node[circ] at (43.287, -10.587){};
	\node[circ] at (43.679, -10.329){};
	\node[circ] at (43.729, -10.579){};
	\node[circ] at (43.679, -10.579){};
	\node[circ] at (43.908, -10.458){};
	\node[circ] at (43.729, -10.579){};
	\node[circ] at (43.679, -10.579){};
	\node[circ] at (43.958, -10.608){};
	\node[circ] at (44.208, -10.608){};
	\node[circ] at (43.908, -10.458){};
	\node[circ] at (44.058, -10.508){};
	\node[circ] at (43.558, -10.508){};
	\node[circ] at (43.108, -11.358){};
	\node[circ] at (44.158, -9.408){};
	\node[circ] at (44.158, -9.408){};
	\node[circ] at (44.158, -9.408){};
	\node[circ] at (41.5, -10.75){};
	\node[circ] at (41, -8.25){};
	\node[circ] at (44.25, -11.25){};
	\node[circ] at (42.75, -9.25){};
	\node[circ] at (42.604, -9.354){};
	\node[circ] at (42.554, -9.354){};
	\node[circ] at (42.604, -9.354){};
	\node[circ] at (42.554, -9.354){};
	\node[circ] at (42.833, -9.383){};
	\node[circ] at (42.933, -9.283){};
	\node[circ] at (42.433, -9.283){};
	\node[circ] at (42.412, -9.613){};
	\node[circ] at (42.412, -9.613){};
	\node[circ] at (42.562, -9.663){};
	\node[circ] at (42.512, -9.613){};
	\node[circ] at (42.662, -9.613){};
	\node[circ] at (42.613, -9.463){};
	\node[circ] at (42.463, -9.613){};
	\node[circ] at (42.463, -9.363){};
	\node[circ] at (42.662, -9.363){};
	\node[circ] at (42.342, -9.542){};
	\node[circ] at (42.742, -9.592){};
	\node[circ] at (42.421, -9.521){};
	\node[circ] at (42.371, -9.371){};
	\node[circ] at (42.421, -9.271){};
	\node[circ] at (42.5, -9.5){};
	\node[circ] at (42.942, -9.492){};
	\node[circ] at (42.892, -9.492){};
	\node[circ] at (42.942, -9.492){};
	\node[circ] at (42.892, -9.492){};
	\node[circ] at (42.771, -9.421){};
	\node[circ] at (42.742, -9.642){};
	\node[circ] at (42.692, -9.642){};
	\node[circ] at (42.742, -9.642){};
	\node[circ] at (42.692, -9.642){};
	\node[circ] at (42.971, -9.671){};
	\node[circ] at (43.071, -9.571){};
	\node[circ] at (42.571, -9.571){};
	\node[circ] at (42.55, -9.9){};
	\node[circ] at (42.55, -9.9){};
	\node[circ] at (42.7, -9.95){};
	\node[circ] at (42.65, -9.9){};
	\node[circ] at (42.8, -9.9){};
	\node[circ] at (42.75, -9.75){};
	\node[circ] at (42.6, -9.9){};
	\node[circ] at (42.6, -9.65){};
	\node[circ] at (42.8, -9.65){};
	\node[circ] at (42.479, -9.829){};
	\node[circ] at (42.879, -9.879){};
	\node[circ] at (42.558, -9.808){};
	\node[circ] at (42.508, -9.658){};
	\node[circ] at (42.558, -9.558){};
	\node[circ] at (42.637, -9.787){};
	\node[circ] at (43.079, -9.779){};
	\node[circ] at (43.029, -9.779){};
	\node[circ] at (43.079, -9.779){};
	\node[circ] at (43.029, -9.779){};
	\node[circ] at (42.908, -9.708){};
	\node[circ] at (43.038, -10.187){};
	\node[circ] at (40.8, -10.9){};
	\node[circ] at (40.75, -10.75){};
	\node[circ] at (40.8, -10.65){};
	\node[circ] at (41, -10.7){};
	\node[circ] at (40.8, -10.9){};
	\node[circ] at (40.75, -10.75){};
	\node[circ] at (40.8, -10.65){};
	\node[circ] at (41, -10.7){};
	\node[circ] at (41.2, -7.8){};
	\node[circ] at (41.25, -7.75){};
	\node[circ] at (41.35, -7.65){};
\end{tikzpicture}
	
\end{center}
\caption{Series of two variables with a correlation coefficient between their empirical means equal to zero and equal to 0.94}	
\label{figure5}
\end{figure}


\subsubsection{Estimation of the confidence level}


In many measurement applications, in addition to the combined standard uncertainty ($u_c(Y)$), it is advisable to report the expanded uncertainty ($U(y) = k_v \cdot u_c(Y)$). 
This allows the definition of a confidence interval with bounds $(Y_0 - k_v \cdot u_c(Y))$ and $(Y_0 + k_v \cdot u_c(Y))$, which has probability $(p)$ (confidence level) of containing the expected value of the measured quantity $(Y)$. 

The coverage factor $(k_p)$ corresponding to the confidence level $(p)$ depends on the probability density function ($p.d.f.$) of the random variable associated with the measurement.\\

In most applications, it can be reasonably assumed that the measured random variable follows a normal distribution. 
Under this assumption, the relationship between confidence level and coverage factor can be obtained from the Gaussian probability density function or, more simply, from tables that directly provide this correspondence. Typical values of $(k_v)$ used in practice are 2 and 3, which correspond approximately to confidence levels of 95.4\% and 99.7\%, respectively.\\

If there are reasonable doubts about the normality of the probability distribution of the measured random variable, Chebyshev’s inequality may be used. This provides an upper bound on the probability that the measurand lies outside the interval $(Y_0 \pm k_v \cdot u_c(Y))$, independently of the form of the probability distribution function. 

However, this upper bound, equal to $(k^{-2})$, is almost always much larger than the actual probability. 
For example, in the case of a Gaussian distribution, the probability that the measurand falls outside the interval $(Y_0 \pm 3 \cdot u_c(Y))$ is approximately 0.003, whereas the upper bound given by Chebyshev’s inequality is about 0.11.\\


A better correspondence between confidence level and coverage factor can be obtained using other techniques not proposed here such as, correlation between the confidence level and the coverage factor that will not be covered in this text.

It should be noted, however, that these techniques are used almost exclusively in high-level measurement applications, i.e., when uncertainties are close to those of primary reference standards.

\subsubsection{Examples}

As a conclusion to this section, and as an illustrative example of the rules presented so far, we report the evaluation of the standard uncertainty of a quantity obtained through a repeated-measurement method.\\

The output voltage of a stabilized power supply is measured using a digital multimeter by applying a repeated-observation method, obtaining 20 measurements of the measurand (example given in the left side of \autoref{figure2728}) \footnote{For the calculation an example with 20 samples was chosen, in the figure only some samples are shown.}. The estimate of the quantity $(V)$ is given by the sample mean $(\overline{V})$, equal to 9.9992 V. By examining the available readings, the observed phenomenon can be represented by the following model:

\begin{equation}
	V_i = \overline{V} + n_i
\label{eq:251}	
\end{equation}

where $n_i$ are the successive realizations of the measurand signal superimposed on the measurement signal.\\

Applying the Type A uncertainty evaluation method, the experimental standard deviation of the noise is estimated as:

\begin{equation}
	u_A(V) = s(n) = 9 \unit{mV}
\end{equation}

Since the estimate of the measurand is given by the empirical mean, the standard uncertainty is equal to the experimental standard deviation of the mean (\autoref{eq:237}), evaluated as:

\begin{equation}
	u_A(V) = \frac{s(n)}{\sqrt{20}} \approx 2 \text{ mV} 
\end{equation}

In addition to the contribution related to the repeatability of the measurand, one must also account for the fact that each reading $V_i$ is affected by instrumental uncertainty. 

Assume that the uncertainty specifications of the instrument are expressed as a value interval of width $2 \delta_v$, with $\delta_v = 2.5 $ mV.

Under the assumption that all values within this interval are equally likely to represent the measurand i.e., assuming a uniform probability density function for the random variable the standard uncertainty of each observation evaluated by a Type B method is given by:

\begin{equation}
	u_B(V_i) = \frac{\delta_v}{\sqrt{3}} \approx 1.44 \text{ mV}, \quad i = 1, \ldots, 20
\end{equation}

The propagation of the Type B uncertainty of the individual observations onto the estimate $V$ of the measurand can be carried out by noting that the sample mean estimator (\autoref{eq:233}) can be treated as an indirect measurement method, so that the rules presented in \autoref{uncprop243} can be applied.\\

\begin{figure}[h!]
\begin{center}

\begin{tikzpicture}[transform shape]
	% Paths, nodes and wires:
	\draw[line width=0.61234pt] (33, -7) -- (37, -7);
	\draw[line width=0.61234pt] (33, -11) -- (37, -11);
	\draw[line width=0.61234pt] (33, -7) -- (33, -11);
	\draw[line width=0.61234pt] (37, -7) -- (37, -11);
	\node[shape=rectangle, minimum width=0.965cm, minimum height=0.965cm] at (35, -6.25){} node[anchor=center, align=center, text width=0.577cm, inner sep=6pt] at (35, -6.25){\large $V_1(V)$};
	\node[shape=rectangle, minimum width=0.965cm, minimum height=0.965cm] at (41, -6.25){} node[anchor=center, align=center, text width=0.577cm, inner sep=6pt] at (41, -6.25){\large $V_1(V)$};
	\node[shape=rectangle, minimum width=0.465cm, minimum height=0.465cm] at (32, -7){} node[anchor=center, align=center, text width=0.077cm, inner sep=6pt] at (32, -7){10.02};
	\node[shape=rectangle, minimum width=0.465cm, minimum height=0.465cm] at (32, -8){} node[anchor=center, align=center, text width=0.077cm, inner sep=6pt] at (32, -8){10.01\\};
	\node[shape=rectangle, minimum width=0.465cm, minimum height=0.465cm] at (32, -9){} node[anchor=center, align=center, text width=0.077cm, inner sep=6pt] at (32, -9){10.00\\};
	\node[shape=rectangle, minimum width=0.465cm, minimum height=0.465cm] at (32, -10){} node[anchor=center, align=center, text width=0.077cm, inner sep=6pt] at (32, -10){9.99\\};
	\node[shape=rectangle, minimum width=0.465cm, minimum height=0.465cm] at (32, -11){} node[anchor=center, align=center, text width=0.077cm, inner sep=6pt] at (32, -11){9.98\\};
	\draw[line width=0.3pt] (33, -9) -- (37, -9);
	\draw[line width=0.3pt, dash pattern={on 0.3pt off 2.4pt}] (33, -8) -- (37, -8);
	\draw[line width=0.3pt, dash pattern={on 0.3pt off 2.4pt}] (33, -10) -- (37, -10);
	\draw[line width=0.3pt, dash pattern={on 0.3pt off 2.4pt}] (34, -7) -- (34, -11);
	\draw[line width=0.3pt, dash pattern={on 0.3pt off 2.4pt}] (35, -7) -- (35, -11);
	\draw[line width=0.3pt, dash pattern={on 0.3pt off 2.4pt}] (36, -7) -- (36, -11);
	\node[shape=rectangle, minimum width=0.465cm, minimum height=0.465cm] at (33, -11.75){} node[anchor=center, align=center, text width=0.077cm, inner sep=6pt] at (33, -11.75){0\\};
	\node[shape=rectangle, minimum width=0.465cm, minimum height=0.465cm] at (35, -11.75){} node[anchor=center, align=center, text width=0.077cm, inner sep=6pt] at (35, -11.75){10\\};
	\node[shape=rectangle, minimum width=0.465cm, minimum height=0.465cm] at (34, -11.75){} node[anchor=center, align=center, text width=0.077cm, inner sep=6pt] at (34, -11.75){5\\};
	\node[shape=rectangle, minimum width=0.465cm, minimum height=0.465cm] at (36, -11.75){} node[anchor=center, align=center, text width=0.077cm, inner sep=6pt] at (36, -11.75){15\\};
	\node[shape=rectangle, minimum width=0.465cm, minimum height=0.465cm] at (37, -11.75){} node[anchor=center, align=center, text width=0.077cm, inner sep=6pt] at (37, -11.75){20};
	\node[shape=rectangle, minimum width=0.465cm, minimum height=0.465cm] at (39, -11.75){} node[anchor=center, align=center, text width=0.077cm, inner sep=6pt] at (39, -11.75){0\\};
	\node[shape=rectangle, minimum width=0.465cm, minimum height=0.465cm] at (41, -11.75){} node[anchor=center, align=center, text width=0.077cm, inner sep=6pt] at (41, -11.75){10\\};
	\node[shape=rectangle, minimum width=0.465cm, minimum height=0.465cm] at (40, -11.75){} node[anchor=center, align=center, text width=0.077cm, inner sep=6pt] at (40, -11.75){5\\};
	\node[shape=rectangle, minimum width=0.465cm, minimum height=0.465cm] at (42, -11.75){} node[anchor=center, align=center, text width=0.077cm, inner sep=6pt] at (42, -11.75){15\\};
	\node[shape=rectangle, minimum width=0.465cm, minimum height=0.465cm] at (43, -11.75){} node[anchor=center, align=center, text width=0.077cm, inner sep=6pt] at (43, -11.75){20};
	\node[circ] at (33.25, -9.75){};
	\node[circ] at (33.75, -7.75){};
	\node[circ] at (34.25, -8.75){};
	\node[circ] at (35.5, -10){};
	\node[circ] at (34.489, -10.25){};
	\node[circ] at (35, -8.25){};
	\node[circ] at (35.25, -9.5){};
	\node[circ] at (36.5, -9.75){};
	\node[circ] at (36.75, -7.75){};
	\draw[line width=0.61234pt] (39, -7) -- (43, -7);
	\draw[line width=0.61234pt] (39, -7) -- (39, -11);
	\draw[line width=0.61234pt] (43, -7) -- (43, -11);
	\draw[line width=0.3pt] (39, -9) -- (43, -9);
	\draw[line width=0.3pt, dash pattern={on 0.3pt off 2.4pt}] (39, -8) -- (43, -8);
	\draw[line width=0.3pt, dash pattern={on 0.3pt off 2.4pt}] (39, -10) -- (43, -10);
	\draw[line width=0.3pt, dash pattern={on 0.3pt off 2.4pt}] (40, -7) -- (40, -11);
	\draw[line width=0.3pt, dash pattern={on 0.3pt off 2.4pt}] (41, -7) -- (41, -11);
	\draw[line width=0.3pt, dash pattern={on 0.3pt off 2.4pt}] (42, -7) -- (42, -11);
	\node[circ] at (39.25, -9.75){};
	\node[circ] at (39.75, -7.75){};
	\node[circ] at (40.25, -8.75){};
	\node[circ] at (41.5, -10){};
	\node[circ] at (40.489, -10.25){};
	\node[circ] at (41, -8.25){};
	\node[circ] at (41.25, -9.5){};
	\node[circ] at (42.5, -9.75){};
	\node[circ] at (42.75, -7.75){};
	\draw[line width=0.61234pt] (39, -11) -- (43, -11);
	\draw (42.5, -9.5) -- (42.5, -10) -- (42.6, -10) -- (42.4, -10);
	\draw (42.4, -9.5) -- (42.6, -9.5);
	\draw (39.75, -7.5) -- (39.75, -8) -- (39.85, -8) -- (39.65, -8);
	\draw (39.65, -7.5) -- (39.85, -7.5);
	\draw (41, -8) -- (41, -8.5) -- (41.1, -8.5) -- (40.9, -8.5);
	\draw (40.9, -8) -- (41.1, -8);
	\draw (42.75, -7.5) -- (42.75, -8) -- (42.85, -8) -- (42.65, -8);
	\draw (42.65, -7.5) -- (42.85, -7.5);
	\draw (40.25, -8.5) -- (40.25, -9) -- (40.35, -9) -- (40.15, -9);
	\draw (40.15, -8.5) -- (40.35, -8.5);
	\draw (41.5, -9.75) -- (41.5, -10.25) -- (41.6, -10.25) -- (41.4, -10.25);
	\draw (41.4, -9.75) -- (41.6, -9.75);
	\draw (41.5, -9.75) -- (41.5, -10.25) -- (41.6, -10.25) -- (41.4, -10.25);
	\draw (41.4, -9.75) -- (41.6, -9.75);
	\draw (41.5, -9.75) -- (41.5, -10.25) -- (41.6, -10.25) -- (41.4, -10.25);
	\draw (41.4, -9.75) -- (41.6, -9.75);
	\draw (39.25, -9.5) -- (39.25, -10) -- (39.35, -10) -- (39.15, -10);
	\draw (39.15, -9.5) -- (39.35, -9.5);
	\draw (40.5, -10) -- (40.5, -10.5) -- (40.6, -10.5) -- (40.4, -10.5);
	\draw (40.4, -10) -- (40.6, -10);
	\draw (40.5, -10) -- (40.5, -10.5) -- (40.6, -10.5) -- (40.4, -10.5);
	\draw (40.4, -10) -- (40.6, -10);
	\draw (40.5, -10) -- (40.5, -10.5) -- (40.6, -10.5) -- (40.4, -10.5);
	\draw (40.4, -10) -- (40.6, -10);
	\draw (41.25, -9.25) -- (41.25, -9.75) -- (41.35, -9.75) -- (41.15, -9.75);
	\draw (41.15, -9.25) -- (41.35, -9.25);
\end{tikzpicture}
	
\end{center}
\caption{Previously used voltages readings with and without assigned values band given by the multimeter reading.}
\label{figure2728}
\end{figure}

To estimate the correlation coefficient between the different readings, we consider model (\ref{eq:251}) in the absence of noise $(n_i = 0)$, since its effect has already been accounted for using the Type A method. 

Under these conditions, the correlation coefficient between the observations $V_i$ can be considered equal to unity, since the different readings are obtained with the same instrument and the same measurement range. 

Therefore, relation (\ref{eq:248}) can be applied, yielding:

\begin{equation}
	u_B(\overline{V}) = \sum_{i=1}^{20} \frac{\partial \overline{V}}{\partial V_i} \cdot u_B(V_i) = \frac{1}{20} \cdot 20 \cdot u_B(V_i) = u_B(V_i)
\end{equation}

In the end we get that:

\begin{equation}
	u_B(\overline{V}) = u_B(V_i) = 1.44 \unit{mV}
\end{equation}

which is a reasonable result, as it shows that the averaging operation does not reduce a systematic effect such as instrumental uncertainty.\\

Finally, by combining the contributions obtained through Type A and Type B methods, the standard uncertainty of the sample mean $\overline{V}$ is obtained:

\begin{equation}
	u(\overline{V}) = \sqrt{u^2_A(\overline{V}) + u^2_B(\overline{V})}  \approx 2.55 \unit{mV}
\end{equation}







\end{document}








